\documentclass[11pt,a4paper]{article}

\usepackage[utf8]{inputenc}
\usepackage[T1]{fontenc}
\usepackage{lmodern}
\usepackage{natbib}
\usepackage[french]{babel}
\usepackage[margin=2.5cm]{geometry}
\usepackage{amsmath,amssymb}
\usepackage{graphicx}
\usepackage{booktabs}
\usepackage{longtable}
\usepackage{tabularx}
\usepackage{ragged2e}
\usepackage{threeparttable}
\usepackage{caption}
\usepackage{setspace}
\usepackage{xcolor}
\usepackage{tikz}
\usetikzlibrary{arrows.meta,positioning,fit}
\usepackage[colorlinks=true, linkcolor=blue!60!black, citecolor=blue!60!black, urlcolor=blue!60!black]{hyperref}

\graphicspath{{../../07_reports/figures/}}

\captionsetup{font=small, labelfont=up}
\addto\captionsfrench{\renewcommand{\tablename}{Tableau}}
\setlength{\parskip}{0.2em}
\onehalfspacing

\title{\textbf{Fiscalité, informalité, transferts et pauvreté en Côte d'Ivoire :\\
une analyse d'incidence sur l'EHCVM 2021}\thanks{%
Document de travail de la Cellule d'Analyses Économiques (ANStat). Les sources,
les hypothèses et la procédure de réplication sont documentées en
annexe~\ref{annexe:pipeline}. Les erreurs éventuelles restent les nôtres.}}

\author{Franck MIGONE, Djaha Armand et Lebe N'Dre\\
\normalsize ANStat}

\date{Juillet 2026}

\begin{document}

\maketitle

\begin{abstract}
\noindent Cette étude mesure l'incidence des impôts et des dépenses publiques
en Côte d'Ivoire à partir de l'EHCVM 2021 et du cadre \emph{Commitment to
Equity}. Elle combine les barèmes fiscaux, les budgets exécutés et le tableau
national des ressources et des emplois, en tenant compte des achats informels.
Dans le scénario central, le Gini passe de 0,3414 pour le
revenu primaire à 0,3168 pour le revenu final. Mais la pauvreté monétaire atteint
42,0~\% au revenu consommable, contre 37,7~\% au revenu primaire, et 4,34~\% de
la population passe sous le seuil national. Les services publics ramènent
l'indicateur du revenu final à 33,9~\%, mais ne procurent pas de liquidités aux
ménages. Enfin, une taxe de 9~\% sur la vente finale de certains produits, simulée à
TVA incorporée constante, accroît la pauvreté de 0,47 point sans compensation
et la réduit de 0,23 point lorsque la recette est redistribuée.
\end{abstract}

\vspace{0.5em}
\noindent\textbf{Mots-clés :} incidence fiscale, TVA, transferts, dépenses
publiques, informalité, pauvreté, microsimulation, Côte d'Ivoire.\\
\textbf{Codes JEL :} H22, D31, I32, O17.

\newpage

%=============================================================================
\section{Introduction}
\label{sec:intro}

En Côte d'Ivoire, l'action publique réduit les écarts de niveau de vie sans
protéger tous les ménages contre la pauvreté monétaire. Dans le scénario central
de cette étude, le coefficient de Gini passe de 0,3414 pour le revenu primaire à
0,3168 lorsque les services publics d'éducation et de santé sont ajoutés. Avant
cette valorisation en nature, le taux de pauvreté atteint pourtant 42,0~\% au
revenu consommable, contre 37,7~\% au revenu primaire. Cette tension est au cœur
de l'analyse : un système fiscal peut réduire l'inégalité relative tout en
abaissant les ressources monétaires de ménages proches du seuil de pauvreté.

Nous mesurons l'incidence des impôts et des dépenses publiques en suivant, pour
chaque ménage, la séquence de revenus du cadre \emph{Commitment to Equity} (CEQ)
\citep{lustig2022}. L'analyse combine l'EHCVM 2021, les barèmes fiscaux, les
budgets exécutés, le Tarif extérieur commun de la CEDEAO et le tableau ivoirien
des ressources et des emplois de 2023. Elle porte sur 12~965 ménages,
64~491 personnes et 787~410 lignes de dépenses. Les impôts indirects comprennent
la TVA, les accises et les droits de douane; les dépenses publiques couvrent les
transferts monétaires, les réductions de prix de l'eau et de l'électricité,
ainsi que les services publics d'éducation et de santé. La consommation par tête
de l'enquête sert de point d'entrée pour le revenu disponible, car elle est
mieux mesurée que le revenu courant. Les concepts antérieurs sont reconstruits
à rebours et cette convention est explicitement testée pour le traitement des
pensions.

\subsection*{Repères de lecture}

Les revenus CEQ sont des étapes comptables, non des catégories administratives
de ménages. Chaque passage ajoute un groupe d'instruments afin d'en isoler
l'effet sur le niveau de vie, la pauvreté et l'inégalité.

\begin{table}[h]
\centering
\caption{La logique de lecture de la cascade CEQ}
\label{tab:guide-lecture}
\small
\begin{tabularx}{\textwidth}{>{\RaggedRight\arraybackslash}p{0.22\textwidth}
                            >{\RaggedRight\arraybackslash}p{0.31\textwidth}
                            >{\RaggedRight\arraybackslash}X}
\toprule
Passage & Opérations ajoutées & Question économique \\
\midrule
Revenu primaire $\rightarrow$ revenu disponible
 & Retrancher impôts directs et cotisations; ajouter transferts monétaires
 & De combien l'État modifie-t-il les ressources monétaires directement
   disponibles? \\
Revenu disponible $\rightarrow$ revenu consommable
 & Retrancher TVA, accises et droits de douane; ajouter réductions publiques de prix
 & Quel pouvoir d'achat reste-t-il après les taxes et les aides incorporées
   à la consommation? \\
Revenu consommable $\rightarrow$ revenu final
 & Ajouter la valeur des services publics d'éducation et de santé
 & Que reçoit le ménage lorsque les prestations en nature sont aussi prises
   en compte? \\
\bottomrule
\end{tabularx}
\end{table}

Trois distinctions évitent les principaux contresens. Premièrement, la TVA
\emph{directe} est contenue dans l'achat final. La TVA \emph{non déductible} est
payée sur un intrant sans pouvoir être récupérée par l'entreprise; lorsqu'elle
se transmet de branche en branche jusqu'au prix final, nous la qualifions de TVA
\emph{enchâssée}. Deuxièmement, les conventions PDI et PGT ne séparent pas deux
populations. Elles classent les mêmes pensions contributives soit comme un
revenu différé déjà acquis (PDI), soit comme un transfert public (PGT).
Troisièmement, un impôt est dit progressif lorsque sa charge augmente plus vite
que le revenu de référence. Cette propriété ne garantit ni une baisse de la
pauvreté ni un gain de pouvoir d'achat. Le lexique de la
section~\ref{sec:lexique} définit les termes et les sigles avant leur
formalisation.

La mesure de la TVA distingue le droit fiscal de la collecte effective. Le
scénario central, S3, fait varier la probabilité de taxation selon la fonction de
consommation et le décile du ménage. Une variante par milieu de résidence, S2,
et la transmission intégrale encadrent cette hypothèse. La TVA non déductible est
ensuite propagée par un modèle de prix de Leontief qui sépare intrants locaux et
importés et replace les emplois du TRE aux prix de base. Le TRE 2023 aux prix
courants est la matrice de référence; sa version à prix constants et la base
internationale ICIO 2020 de l'OCDE servent de contrôles.

Les résultats montrent d'abord que l'informalité des circuits de vente modifie
fortement la charge estimée. Sous S3, la TVA directe atteint 693 milliards de
FCFA, contre 1~387 milliards sous transmission intégrale. Le TRE ajoute
214 milliards de TVA enchâssée; les variantes nationale à prix constants et
internationale donnent des masses totales très proches. Les accises et les
droits de douane directement contenus dans les achats ajoutent respectivement
67,9 et 185,1 milliards. Après l'ensemble des interventions monétaires, le
revenu consommable place 42,0~\% de la population sous le seuil national, et
4,34~\% de la population franchit ce seuil. La progressivité relative de plusieurs
prélèvements ne suffit donc pas à protéger les ménages situés près du seuil.

Les prélèvements directs sont plus concentrés, mais leur assiette est étroite.
Le scénario central produit 225,5 milliards de FCFA d'impôt sur le revenu des
personnes physiques et 162,4 milliards de cotisations. Les paiements publics
directs atteignent 125,8 milliards, dont 27,1 milliards pour le Programme de
filets sociaux productifs. Les services publics nets des paiements des familles
représentent 1~301,6 milliards pour l'éducation et 81,85 milliards pour la
santé. Ils expliquent une part importante de la baisse du Gini, mais leur coût de
production n'est pas un revenu liquide. Le taux de pauvreté de 33,9~\% associé
au revenu final est ainsi un indicateur comptable qui ne doit pas être confondu
avec la pauvreté monétaire.

Enfin, l'usage des recettes change le diagnostic d'une réforme de TVA. Une taxe
de 9~\% sur la vente finale de certains produits collecte 97,2 milliards sous
S3 et accroît la pauvreté de 0,47 point si la recette n'est pas restituée. Le
calcul maintient constante la TVA déjà incorporée dans les coûts : il ne simule
pas l'ouverture de nouveaux droits à déduction. Un transfert uniforme de
3~260 FCFA par personne, exactement financé par cette recette, réduit au
contraire la pauvreté de 0,23 point. Ce contrefactuel montre qu'une réforme doit
être évaluée avec l'usage de ses recettes.

\subsection*{Positionnement dans la littérature}

Le premier point d'ancrage est la littérature sur l'incidence fiscale CEQ. Le
manuel dirigé par \citet{lustig2022} organise la comparaison des revenus avant
et après impôts, transferts et services publics. \citet{higgins2016} montrent
qu'un système peut réduire l'inégalité tout en appauvrissant des ménages lorsque
les taxes indirectes dépassent les paiements qu'ils reçoivent. Les applications
au Ghana \citep{younger2016}, à l'Éthiopie \citep{hill2017} et à trois pays
d'Afrique de l'Ouest francophone \citep{akim2020} établissent la pertinence de
ce diagnostic sur le continent. La présente étude actualise le cas ivoirien sur
l'EHCVM 2021 et construit la séquence jusqu'au revenu final, tout en publiant les
conventions de pension, les champs couverts et les écarts avec les agrégats
administratifs.

Le deuxième point d'ancrage concerne l'informalité et les taxes sur la
consommation. \citet{bachas2024} montrent que la part des achats informels tend
à diminuer avec le revenu et que ce profil peut rendre les taxes à la
consommation plus progressives que ne l'indiquent les seuls taux légaux. Les
exonérations créent aussi de la TVA non récupérable qui peut se transmettre par
les chaînes de production \citep{chandler2025}. Nous réunissons ces deux
mécanismes : trois scénarios explicites décrivent la taxation effective des
achats, tandis qu'une concordance relie les postes de consommation aux
48 produits du TRE pour mesurer la charge incorporée aux prix.

Le troisième point d'ancrage compare les exonérations aux transferts. Dans six
pays à revenu faible ou intermédiaire, \citet{warwick2022} montrent que les taux
réduits et les exonérations ciblent imparfaitement les pauvres et qu'un transfert
financé par l'élargissement de l'assiette peut mieux les protéger. La simulation
ivoirienne reprend cette comparaison en tenant compte de l'informalité des
achats. Les intervalles de confiance reposent sur 500 réplications du bootstrap
Rao--Wu \citep{raowu1988}; la recette et le transfert sont réestimés dans chaque
réplication.

La section suivante décrit le système fiscal de 2021. Les sections consacrées
aux données et à la méthode expliquent ensuite les variables, les conventions et
les calculs avant la présentation des résultats. Les simulations de réforme
précèdent l'examen des limites, puis la conclusion tire les implications pour la
conception d'un ensemble cohérent de prélèvements et de paiements publics.

%=============================================================================
\section{Contexte institutionnel : le système fiscal ivoirien en 2021}
\label{sec:contexte}

La structure fiscale ivoirienne repose sur les impôts indirects. La TVA,
instituée au taux normal de 18~\%, coexiste avec un taux réduit de 9~\%
--- appliqué notamment au lait, aux pâtes alimentaires et aux équipements
solaires --- et avec un vaste champ d'exonérations couvrant les produits
alimentaires de base, la santé et l'éducation. La Banque mondiale
\citep{worldbank2019} estime que ces régimes dérogatoires expliquent l'essentiel
de l'écart entre les recettes potentielles et les recettes effectives de TVA, et
en fait le premier gisement de mobilisation fiscale du pays. Cette architecture
a une conséquence directe pour l'analyse d'incidence : le taux moyen pondéré du
panier de consommation est très inférieur au taux légal, et il varie fortement
le long de la distribution des niveaux de vie, puisque la part des produits
exonérés décroît avec le revenu.

Le deuxième trait structurant est l'informalité. L'essentiel de l'emploi et une
grande partie du commerce de détail échappent à l'enregistrement fiscal. Du côté
des achats, cela signifie qu'une fraction importante de la dépense des ménages
--- surtout des ménages pauvres et ruraux --- est réalisée auprès de vendeurs
qui ne collectent pas la TVA ; l'hypothèse de transmission complète du taux
statutaire aux prix payés constitue donc une borne supérieure de la charge.
Du côté des revenus, l'informalité restreint l'assiette des prélèvements
directs : seuls les salariés du secteur formel --- administration, entreprises
enregistrées --- subissent la retenue à la source de l'impôt sur les salaires et
les cotisations sociales.

L'imposition des salaires en 2021 combine trois prélèvements cédulaires
successifs. L'impôt sur les salaires (IS) frappe la base imposable ---
c'est-à-dire 80~\% du salaire brut, après l'abattement forfaitaire de 20~\% pour
frais professionnels --- au taux de 1,5~\%, soit 1,2~\% du brut. La contribution
nationale (CN) applique à la même base un barème progressif par tranches
mensuelles (0~\% jusqu'à 50~000 FCFA, puis 1,5~\%, 5~\% et 10~\% au-delà de
200~000 FCFA). L'impôt général sur le revenu (IGR) s'ajoute enfin selon une
logique de quotient familial : la base, nette d'IS et de CN, est divisée par un
nombre de parts plafonné à cinq (une part pour un célibataire, deux pour un
couple marié, une demi-part par enfant à charge), soumise à un barème mensuel en
dix tranches, puis l'impôt par part est multiplié par le nombre de parts et
réduit de 15~\%. S'y ajoutent la cotisation retraite CNPS à la charge du salarié
(6,3~\% du brut, plafonné à 1~647~315 FCFA par mois) et la contribution à la
Couverture Maladie Universelle (CMU), fixée à 1~000 FCFA par mois et par adulte
contributif participant. La participation est simulée selon les cibles et la
couverture déclarée dans l'enquête : les personnes dont l'assurance maladie est
financée par un organisme public forment le vivier observable, ensuite séparé
entre pauvres non contributifs et adultes non pauvres contributifs. Les cibles,
la règle de tirage et la variante universelle sont détaillées en
section~\ref{sec:methode}. L'annexe~\ref{annexe:baremes} reproduit les barèmes.

La fiscalité indirecte ne se limite pas à la TVA. Le TEC CEDEAO répartit les
produits importés entre les bandes de 0, 5, 10, 20 et 35~\%; les accises 2021
frappent notamment alcools, tabacs, produits pétroliers, boissons non
alcoolisées et certains cosmétiques. Côté dépenses, le PSSN verse des transferts
monétaires ciblés tandis que les bourses et les branches familiales et accidents
du travail de la protection sociale complètent les transferts directs. Ces
instruments sont traités séparément parce que leur assiette, leur mécanisme de
ciblage et leur incidence économique diffèrent.

%=============================================================================
\section{Données}
\label{sec:donnees}

Nous utilisons l'EHCVM 2021 de la Côte d'Ivoire, enquête nationale représentative
conduite dans le cadre harmonisé de l'Union économique et monétaire
ouest-africaine (UEMOA). Le module de consommation fournit
787~410 observations au niveau ménage~$\times$~produit~$\times$~mode
d'acquisition pour 12~965 ménages. Le fichier individuel couvre 64~491 personnes.
Le fichier officiel de bien-être apporte la consommation annuelle par tête, le
déflateur spatial, le seuil de pauvreté officiel, les poids d'enquête, les
grappes et les strates. Le revenu disponible de référence est donc l'agrégat de
consommation par tête dans les prix réels de l'enquête; les charges fiscales,
initialement calculées en FCFA courants, sont converties avec le déflateur
spatial avant d'en être soustraites.

Le questionnaire est mobilisé au-delà du fichier harmonisé de bien-être. Le
tableau~\ref{tab:variables-questionnaire} présente les informations par leur
libellé et précise leur fonction dans la simulation. Il permet de distinguer
les variables effectivement utilisées des informations déjà repérées pour les
étapes futures, sans imposer au lecteur les noms techniques de la base.

\begin{table}[t]
\centering
\begin{threeparttable}
\caption{Informations du questionnaire mobilisées dans l'analyse}
\label{tab:variables-questionnaire}
\small
\begin{tabularx}{\textwidth}{>{\raggedright\arraybackslash}p{3cm}
  >{\raggedright\arraybackslash}X >{\raggedright\arraybackslash}X}
\toprule
Module & Informations observées & Utilisation dans le papier \\
\midrule
Niveau de vie & Consommation annuelle par tête, déflateur spatial, seuil de
pauvreté, poids, strate et grappe & Ancrage de $Y_D$, classement des ménages,
prix réels, pauvreté et inférence de sondage \\
Consommation & Produit consommé, valeur annuelle et mode d'acquisition &
Assiette de TVA, taux statutaire, scénario de taxation effective et
concordance avec les secteurs du TRE \\
Emploi principal & Nombre de mois travaillés, salaire et périodicité,
bulletin de salaire, cotisation à la Caisse générale de retraite des agents de
l'État (CGRAE) ou à la CNPS, congés payés, jours de
congé annuel, congé maladie et congé de maternité ou paternité & Annualisation
du salaire et construction du scénario central d'assujettissement
multi-critères \\
Emploi secondaire & Existence de l'emploi, profession, catégorie d'activité,
salaire et périodicité & Complément de salaire et identification des emplois
secondaires susceptibles de faire l'objet d'une retenue \\
Santé et assurance & Couverture par une assurance maladie et organisme qui la
finance & Vivier observable de la CMU et distinction contributive ou non
contributive \\
Revenus et programmes & Pensions, revenus fonciers, dividendes, intérêts,
bourses, transferts publics, filets sociaux et transferts privés & Pensions,
PSSN, bourses, prestations familiales, accidents du travail et scénarios PDI/PGT \\
Logement et énergie & Dépenses d'eau, d'électricité et de carburants; type
d'abonnement électrique & Attribution des réductions de prix et robustesses \\
Éducation et santé & Scolarisation, type d'établissement, recours aux soins,
prestataire et dépenses acquittées & Attribution des dépenses publiques,
paiements directs et bénéfices nets \\
\bottomrule
\end{tabularx}
\begin{tablenotes}\footnotesize
\item \emph{Notes} : les transferts privés reçus sont conservés séparément des
transferts publics non contributifs. Les codes informatiques et la
correspondance exacte avec les fichiers sources relèvent du matériel de
réplication, non de la définition économique des variables.
\end{tablenotes}
\end{threeparttable}
\end{table}

Cet inventaire détermine les sources mobilisées dans la méthode. Les pensions et
les paiements publics séparent les concepts de revenu monétaire; les modules de
logement, d'éducation et de santé identifient les ménages auxquels sont attribués
les aides de prix et les services publics, complétés par les paramètres externes
documentés ci-dessous.
Les envois reçus des proches restent des transferts privés et ne sont pas
comptés parmi les transferts publics non contributifs.

Le tableau~\ref{tab:descriptif} donne les principales dimensions de
l'échantillon après jointure. Les pourcentages de salariés et d'assujettis sont
calculés sur l'ensemble de la population, enfants compris; ils décrivent donc la
portée démographique des prélèvements, non le taux de formalité de la population
active. Les poids individuels sont les poids ménage appliqués à chaque membre;
les statistiques de ménages conservent le poids ménage sans multiplication par
la taille.

\begin{table}[t]
\centering
\begin{threeparttable}
\caption{Statistiques descriptives de l'échantillon EHCVM 2021}
\label{tab:descriptif}
\small
\begin{tabular}{lrrrrrr}
\toprule
Zone & Ménages & Personnes & Taille du & Âge & Moins de & Assujettis \\
 & observés & observées & ménage & moyen & 18 ans (\%) & centraux (\%) \\
\midrule
National & 12~965 & 64~491 & 4,62 & 23,8 & 48,0 & 3,61 \\
Urbain   &  5~243 & 25~173 & 4,46 & 24,2 & 44,7 & 5,70 \\
Rural    &  7~722 & 39~318 & 4,82 & 23,3 & 51,6 & 1,29 \\
\bottomrule
\end{tabular}
\begin{tablenotes}\footnotesize
\item \emph{Notes} : la population pondérée est de 30,09 millions de
personnes. La part de salariés rémunérés sur l'ensemble de la population vaut
8,02~\% au niveau national, 11,92~\% en milieu urbain et 3,68~\% en milieu
rural. Source : calculs des auteurs à partir de l'EHCVM 2021.
\end{tablenotes}
\end{threeparttable}
\end{table}

La simulation de TVA retient les dépenses déclarées comme achats et comme
valeurs d'usage. L'agrégat de niveau de vie continue de compter
l'autoconsommation, les dons et les loyers imputés, mais aucune TVA n'est
calculée sur ces éléments puisqu'ils ne correspondent pas à un achat monétaire
taxable. Les dépenses sont winsorisées au 99\textsuperscript{e}
centile à l'intérieur de chaque produit; les versions au 95\textsuperscript{e}
et au 99,5\textsuperscript{e} centile sont conservées pour la sensibilité. Le
traitement par produit évite qu'un seuil dominé par les postes coûteux écrase la
variation légitime des petits postes. Enfin, chaque poste de consommation est
apparié à son taux statutaire de 0, 9 ou 18~\%. Le même poste est relié aux
secteurs entrées--sorties (I/O) par la concordance sectorielle décrite ci-dessous; aucun taux
sectoriel n'est obtenu par une moyenne non pondérée de correspondances multiples.

%=============================================================================
\section{Méthodologie}
\label{sec:methode}

\subsection{Lexique et conventions de lecture}
\label{sec:lexique}

Le tableau~\ref{tab:lexique} donne le sens précis des termes employés dans la
suite. Il distingue notamment une taxe prévue par le droit, la taxe
effectivement contenue dans un prix et la charge économique finalement imputée
au ménage.

\begin{longtable}{>{\RaggedRight\arraybackslash}p{0.24\textwidth}
                  >{\RaggedRight\arraybackslash}p{0.70\textwidth}}
\caption{Lexique minimal de l'analyse}\label{tab:lexique}\\
\toprule
Terme & Définition utilisée dans le document \\
\midrule
\endfirsthead
\toprule
Terme & Définition utilisée dans le document \\
\midrule
\endhead
\bottomrule
\endfoot

\textbf{Incidence fiscale}
& Attribution de la charge économique d'un impôt, ou du bénéfice d'une dépense
publique, aux ménages qui la supportent ou en bénéficient. Elle peut différer
du contribuable qui verse juridiquement l'impôt à l'administration. \\[0.35em]

\textbf{Assiette, taux statutaire et taux effectif}
& L'assiette est le montant auquel le droit applique l'impôt. Le taux
statutaire est celui du barème légal. Le taux effectif rapporte la charge
réellement simulée au revenu ou à la dépense du ménage; il tient donc compte des
exonérations, de l'informalité et de la composition du panier. \\[0.35em]

\textbf{TVA directe}
& Part de taxe sur la valeur ajoutée directement contenue dans le prix toutes
taxes comprises (TTC) des achats finaux du ménage. Pour un prix TTC et un taux
$\tau$, cette part est $\tau/(1+\tau)$, modulée ici par la probabilité que la
transaction soit taxée et que la taxe soit transmise au prix. \\[0.35em]

\textbf{TVA non déductible}
& TVA payée par une entreprise sur ses intrants mais qu'elle ne peut pas
récupérer comme crédit de TVA, notamment lorsque sa production est exonérée ou
vendue hors du circuit déclaratif. Elle devient alors un coût de production. \\[0.35em]

\textbf{TVA enchâssée}
& Part cumulée de TVA non déductible qui se propage d'un fournisseur à l'autre
et atteint finalement le prix d'un bien acheté par le ménage. Il ne s'agit pas
d'un impôt juridique supplémentaire, mais d'une charge incorporée au prix par
les chaînes de production. \\[0.35em]

\textbf{Accise}
& Impôt indirect ciblé sur certains produits, par exemple carburants, tabacs,
alcools, boissons non alcoolisées ou cosmétiques. Il peut être proportionnel à
la valeur ou fixé par unité physique et s'ajoute, selon la cascade légale, aux
droits de douane et à la TVA. \\[0.35em]

\textbf{Droit de douane}
& Impôt acquitté lors de l'importation. La charge imputée au ménage dépend du
taux du Tarif extérieur commun de la CEDEAO et de la part importée du produit
consommé. \\[0.35em]

\textbf{TRE et modèle de Leontief}
& Le tableau des ressources et des emplois (TRE) décrit, dans les comptes
nationaux, quels produits sont fournis et utilisés par chaque branche. Le
modèle de prix de Leontief utilise ces relations pour suivre la transmission
d'un coût d'intrant vers les branches en aval, à quantités et technologies
fixes. \\[0.35em]

\textbf{EHCVM et COICOP}
& L'EHCVM est l'Enquête harmonisée sur les conditions de vie des ménages. La
COICOP est la nomenclature internationale qui classe leur consommation par
fonction, par exemple alimentation, logement ou transport. \\[0.35em]

\textbf{CGI, TEC et CEDEAO}
& Le Code général des impôts (CGI) fixe les barèmes nationaux. Le Tarif
extérieur commun (TEC) répartit les importations entre bandes tarifaires
communes aux pays de la Communauté économique des États de l'Afrique de l'Ouest
(CEDEAO). \\[0.35em]

\textbf{ICIO et analyse entrées--sorties}
& ICIO désigne la base internationale de tableaux entrées--sorties de l'OCDE,
utilisée ici comme robustesse au TRE national. Une analyse entrées--sorties
décrit les achats d'intrants entre branches et la transmission des coûts qui en
résulte. \\[0.35em]

\textbf{CEQ}
& \emph{Commitment to Equity}: cadre d'analyse qui compare plusieurs concepts
de revenu avant et après impôts, cotisations, transferts monétaires,
subventions et prestations publiques en nature. \\[0.35em]

\textbf{Concepts de revenu}
& $Y_P$ est le revenu primaire reconstruit avant prélèvements directs et
transferts; $Y_N=Y_P-D-C$ est le revenu net de marché; $Y_G=Y_P+R$ est le revenu
brut après transferts mais avant prélèvements; $Y_D=Y_N+R$ est le revenu
disponible; $Y_C=Y_D-I+S$ est le revenu consommable après impôts et subventions
indirects; $Y_F$ ajoute les prestations publiques en nature. $Y_N$ et $Y_G$
sont deux branches comptables parallèles, non deux étapes successives. \\[0.35em]

\textbf{PDI}
& \emph{Pensions as Deferred Income}, soit « pensions comme revenu différé ».
Les pensions contributives appartiennent au revenu avant intervention et les
cotisations retraite sont traitées comme de l'épargne. Les pensions ne sont
donc pas comptées parmi les transferts publics directs. \\[0.35em]

\textbf{PGT}
& \emph{Pensions as Government Transfer}, soit « pensions comme transfert
public ». Les pensions sont ajoutées aux transferts directs et les cotisations
retraite aux prélèvements. PDI et PGT sont deux conventions comptables
alternatives appliquées aux mêmes ménages, pas deux populations. \\[0.35em]

\textbf{IRPP, CNPS, CGRAE et CMU}
& L'IRPP désigne ici l'impôt simulé sur les revenus des personnes physiques. La
CNPS est la Caisse nationale de prévoyance sociale, la CGRAE la Caisse générale
de retraite des agents de l'État et la CMU la Couverture maladie universelle.
Les cotisations ne sont pas des impôts au sens juridique,
mais elles sont retranchées avec les prélèvements directs dans la cascade CEQ. \\[0.35em]

\textbf{PSSN et PMT}
& Le Programme de filets sociaux productifs (PSSN) verse des transferts
monétaires aux ménages ciblés. Le \emph{proxy means test} (PMT), ou test
indirect de niveau de vie, classe les ménages à partir de caractéristiques
observables lorsque le registre administratif n'est pas disponible. \\[0.35em]

\textbf{Calage}
& Ajustement explicite d'un effectif, d'un poids ou d'un montant simulé pour
retrouver une référence administrative. Le document précise toujours ce qui est
ajusté : pour le PSSN, le classement des ménages reste inchangé et seul le
montant est multiplié par un facteur commun. \\[0.35em]

\textbf{Réduction publique de prix}
& Écart entre le prix payé par le ménage et un prix de référence lorsque la
politique publique ou le barème social réduit la facture. Cette aide peut être
budgétaire, comme l'aide d'exploitation électrique, ou financée par une
subvention croisée entre catégories d'usagers, comme la tranche sociale de
l'eau. \\[0.35em]

\textbf{Bénéfice en nature, brut et net}
& Coût public moyen d'un service d'éducation ou de santé attribué à son
utilisateur. Le bénéfice brut est le coût attribué; le bénéfice net retranche
les frais directement payés par le ménage. Aucun des deux n'est de l'argent
versé au ménage. \\[0.35em]

\textbf{Valeur de Shapley}
& Part de la variation d'un indicateur attribuée à un instrument après avoir
moyenné sa contribution dans tous les ordres possibles d'introduction. Les
parts s'additionnent exactement à la variation totale. \\[0.35em]

\textbf{Appauvrissement fiscal}
& Perte de ressources monétaires sous le seuil de pauvreté provoquée par les
prélèvements nets des paiements publics. Un ménage déjà pauvre peut être
appauvri sans devenir un « nouveau pauvre »; les deux indicateurs sont donc
publiés séparément. \\[0.35em]

\textbf{Décile, masse et pondération}
& Les déciles partagent la population, classée du niveau de vie le plus faible
au plus élevé, en dix groupes de même taille pondérée. La « masse » d'un impôt
ou transfert est sa somme nationale estimée après application des poids
d'enquête; un ménage pondéré représente le nombre de ménages de la population
auquel répond l'observation enquêtée. \\[0.35em]

\textbf{Scénario central et robustesse}
& Le scénario central est l'ensemble d'hypothèses privilégié pour présenter les
résultats. Une robustesse recalcule le même indicateur sous une hypothèse
alternative plausible. S2 et S3 sont ainsi deux hypothèses paramétriques sur la taxation effective,
et PDI/PGT deux conventions de pension. \\[0.35em]

\textbf{Gini et concentration}
& Le coefficient de Gini mesure l'inégalité d'un revenu entre zéro, égalité
parfaite, et un, concentration maximale. Un indice de concentration applique
une logique comparable à un impôt ou un transfert, mais conserve le classement
des ménages selon leur revenu de référence. \\[0.35em]

\textbf{Kakwani}
& Indice de progressivité d'un prélèvement, égal à la concentration de sa charge
moins le Gini du revenu de classement. Une valeur positive signifie que la
charge est relativement plus concentrée vers le haut de la distribution; elle
ne dit pas, à elle seule, si la pauvreté augmente ou diminue. \\[0.35em]

\textbf{Reynolds--Smolensky}
& Variation du Gini entre le revenu après et le revenu avant intervention. Le
papier utilise $RS=G_{\text{après}}-G_{\text{avant}}$ : une valeur négative
indique donc une réduction de l'inégalité et une valeur positive une
augmentation. \\[0.35em]

\textbf{FGT et intervalle de confiance}
& Les indices de Foster--Greer--Thorbecke mesurent le taux de pauvreté
($P_0$), sa profondeur ($P_1$) et sa sévérité ($P_2$). Un intervalle de
confiance à 95~\% quantifie l'incertitude d'échantillonnage; il est obtenu ici
par bootstrap Rao--Wu, qui respecte strates, grappes et poids de l'enquête. \\
\end{longtable}

\subsection{La séquence de revenus CEQ et le point d'entrée ivoirien}

Le papier distingue les concepts de revenu selon la séquence CEQ. Dans les
équations qui suivent, $D$ désigne les impôts directs payés par le ménage, $C$
ses cotisations sociales, $R$ les paiements monétaires qu'il reçoit des
administrations publiques, $I$ les taxes supportées lors de la consommation et
$S$ les aides publiques qui réduisent le prix payé. Enfin, $H$, $E$ et $O$
mesurent la valeur des services publics reçus en santé, en éducation et dans les
autres domaines. La cascade s'écrit
\begin{align}
Y_N &= Y_P-D-C, &
Y_G &= Y_P+R, \label{eq:ceq-branches}\\
Y_D &= Y_N+R, &
Y_C &= Y_D-I+S, \label{eq:ceq-cash}\\
Y_F &= Y_C+H+E+O. \label{eq:ceq-final}
\end{align}
$Y_N$ et $Y_G$ sont deux branches parallèles : le premier déduit les
prélèvements directs avant transferts, le second ajoute les transferts avant
prélèvements. Le revenu imposable administratif n'est pas un concept de
bien-être CEQ et n'est donc pas intercalé entre $Y_G$ et $Y_D$.

L'EHCVM fournit directement l'ancre $Y_D$ par son agrégat annuel de
consommation par tête. Nous reconstruisons les concepts pré-fiscaux à rebours.
Dans le scénario PDI central, la pension contributive est considérée comme la
restitution d'un revenu acquis au cours de la vie active, et non comme une aide
de l'État. La cotisation retraite est alors une épargne obligatoire, et le
revenu de marché coïncide avec le revenu pré-fiscal ($Y_M=Y_P$). Les sommes du
PSSN, les bourses publiques, les prestations familiales et les indemnisations
d'accidents du travail sont, elles, comptées comme des paiements publics reçus
par les ménages. Dans le scénario de robustesse PGT, la pension est au contraire
comptée comme un paiement public et la cotisation retraite comme un prélèvement.
Ces deux conventions aboutissent exactement au même revenu disponible pour
chaque ménage; seule sa décomposition entre ressources initiales, prélèvements
et paiements publics change.

Le calcul complet retranche du revenu disponible la TVA payée sur les achats
finaux, la TVA non récupérable répercutée dans les prix, les accises et les
droits de douane, puis ajoute les réductions de prix de l'eau et de
l'électricité. Il ajoute enfin les services publics d'éducation et de santé.
Toutes les grandeurs $Y_{\cdot h}$ sont annuelles, réelles et exprimées par
personne. Les instruments sont d'abord agrégés au ménage :
\begin{align}
I_h &= \mathrm{TVA}^{\mathrm{directe}}_h+
\mathrm{TVA}^{\mathrm{non\ deduct.}}_h+\mathrm{Accises}_h+\mathrm{Douane}_h,\\
i_h&=I_h/n_h,\qquad s_h=S_h/n_h,\qquad e_h=E_h^{\mathrm{net}}/n_h,
\qquad q_h=H_h^{\mathrm{net}}/n_h,\\
Y_{Ch} &= Y_{Dh}-i_h+s_h,\\
Y_{Fh} &= Y_{Ch}+e_h+q_h.
\label{eq:yc-complet}
\end{align}
Ici, $n_h$ est la taille du ménage; $S_h$, $E_h^{\mathrm{net}}$ et
$H_h^{\mathrm{net}}$ sont respectivement les réductions de prix, le service
d'éducation et le service de santé, tous déflatés spatialement avant la division
par $n_h$. Le scénario central retranche des services les paiements directs
observables. Les équations ne mélangent donc jamais un montant total du ménage
et un revenu par personne.

\begin{figure}[t]
\centering
\begin{tikzpicture}[
  node distance=5mm,
  box/.style={draw, rounded corners=2pt, align=center, text width=8.2cm,
              minimum height=8mm, fill=green!8},
  arr/.style={-{Latex[length=2mm]}, thick}
]
\node[box] (yp) {$Y_P$ : revenu primaire reconstruit\\
\footnotesize conventions PDI et PGT publiées en parallèle};
\node[box, below=of yp] (yn) {$Y_N=Y_P-D-C$ : revenu net de marché\\
\footnotesize impôts directs et cotisations selon la convention de pension};
\node[box, below=of yn] (yd) {$Y_D=Y_N+R$ : revenu disponible\\
\footnotesize ancre observée; paiements publics monétaires intégrés};
\node[box, below=of yd] (yc) {$Y_C=Y_D-I+S$ : revenu consommable\\
\footnotesize TVA, accises, douane, eau et électricité intégrées};
\node[box, below=of yc] (yf) {$Y_F=Y_C+H+E+O$ : revenu final\\
\footnotesize éducation et santé nettes des paiements directs; $O=0$};
\draw[arr] (yp) -- node[right]{\footnotesize $-D-C$} (yn);
\draw[arr] (yn) -- node[right]{\footnotesize $+R$} (yd);
\draw[arr] (yd) -- node[right]{\footnotesize $-I+S$} (yc);
\draw[arr] (yc) -- node[right]{\footnotesize $+H+E+O$} (yf);
\end{tikzpicture}
\caption{Relations entre les concepts de revenu CEQ. Les égalités comptables
sont vérifiées pour chaque ménage.}
\label{fig:ceq-cascade}
\end{figure}
\subsection{TVA directe et probabilité de taxation effective}

La dépense $d_{kh}$ observée pour le produit $k$ est un montant toutes taxes
comprises. Lorsque seule la TVA facturée sur la vente finale est calculée, la
part de taxe vaut $\tau_k/(1+\tau_k)$, et non le taux appliqué directement au
montant TTC :
\begin{equation}
T_h^{d,\,\text{isolée}}=\sum_k d_{kh}\,\alpha_{kh}
\frac{\tau_k}{1+\tau_k}.
\label{eq:vat-direct}
\end{equation}
Cette formule sert aux diagnostics dans lesquels la TVA incorporée est absente.
Les résultats centraux combinent les composantes directe et incorporée; ils
utilisent alors le dénominateur commun de l'équation~\eqref{eq:io-incidence}.

Le coefficient $\alpha_{kh}\in[0,1]$ résume la collecte et la transmission de la
TVA lors de la vente finale. Le scénario strict pose $\alpha_{kh}=1$. S2 fait
varier ce coefficient par fonction COICOP et milieu de résidence. Dans S3, le
ménage est remplacé par son décile $d$ et chaque fonction $k$ suit la règle
\begin{equation}
\alpha_{kd}=\min\{a_k+(d-1)s_k,1\},\qquad d=1,\ldots,10,
\label{eq:alpha-s3}
\end{equation}
où $a_k$ est le niveau du premier décile et $s_k$ la pente par décile. Les
valeurs exactes de $a_k$, $s_k$ et de la matrice $\alpha_{kd}$ figurent dans
l'annexe~\ref{annexe:s3-parameters}. Elles sont imposées avant l'estimation :
elles ne sont ni estimées sur l'EHCVM, ni calées sur une recette administrative
ivoirienne. S3 transpose seulement le profil qualitatif documenté par
\citet{bachas2024}, selon lequel la part des achats formels augmente avec le
niveau de vie. Les scénarios strict et S2 mesurent la sensibilité des résultats
à cette hypothèse structurante.

\subsection{Du TRE aux coefficients techniques}

Le TRE 2023 aux prix courants est la source de référence. Sa matrice d'emplois
$U^a$ est publiée aux prix d'acquisition. Nous retirons, produit par produit, les
marges de commerce, les marges de transport et les impôts nets sur les produits.
Les marges retirées ne disparaissent pas : elles sont réaffectées aux lignes de
services de commerce et de transport des branches utilisatrices.
En notant $m_i^g$, $m_i^t$ et $q_i$ les parts de ces trois éléments dans la
ressource au prix d'acquisition du produit $i$, la composante propre du produit
est
\begin{equation}
U^{b}_{ij}=U^{a}_{ij}(1-m_i^g-m_i^t-q_i),
\label{eq:basic-prices}
\end{equation}
avant réallocation des deux marges. Cette passerelle réduit les emplois
intermédiaires de 35~241 à 33~869 milliards de FCFA et évite de traiter une taxe
ou une marge comme un intrant physique supplémentaire.

Chaque ligne est ensuite séparée entre origine domestique et importée au prorata
de la ressource correspondante. Si $s_i=x_i^d/(x_i^d+m_i)$ est la part de
l'offre domestique, alors $U^d_{ij}=s_iU^b_{ij}$ et
$U^m_{ij}=(1-s_i)U^b_{ij}$. Sous l'hypothèse de technologie de branche, les
matrices de coefficients produits-produits sont
\begin{equation}
B^d=U^d\widehat{x}^{-1},\qquad
B^m=U^m\widehat{x}^{-1},\qquad
T=V'\widehat{g}^{-1},\qquad
A^d=B^dT,\quad A^m=B^mT,
\label{eq:sut-product}
\end{equation}
où $V$ est la matrice de production, $x$ la production des branches, $g$ la
production des produits et $T$ la matrice qui convertit les coefficients par
branche en coefficients par produit. Le rayon spectral de $A^d$ vaut 0,346 pour le TRE
courant, 0,347 pour le TRE constant et 0,318 pour ICIO; ils sont tous inférieurs
à un. Les inverses de Leontief existent et les résidus numériques des équations
sont de l'ordre de $10^{-17}$.

\subsection{Concordance sectorielle et vecteurs fiscaux}

La concordance est construite en deux passages déterministes. Premièrement,
chaque poste de consommation est affecté, selon la nature du bien ou du service,
à une seule catégorie d'une nomenclature intermédiaire de 52 secteurs.
Deuxièmement, chacune de ces catégories est reliée à un seul des 48 produits du
TRE. Un poste ne peut donc être ni fractionné entre plusieurs branches ni
dupliqué. Cette règle conserve la destination sectorielle détaillée de chaque
dépense et évite de répartir un produit entre plusieurs secteurs par une moyenne
non pondérée. La concordance complète figure dans l'annexe~\ref{annexe:mapping-io}. La couverture
pondérée de la consommation est de 92,5~\% pour le TRE; l'ICIO, dont la
concordance est distincte, couvre 95,2~\%.

Pour chaque produit TRE $i$, $\bar\tau_i$ est le taux statutaire moyen appliqué
aux intrants et $e_i$ approche la fraction de production qui ouvre un droit à
déduction. Les deux sont construits à partir de la dépense nationale de
l'EHCVM :
\begin{equation}
\bar\tau_i=
\frac{\sum_{h,k\in i}w_h d_{kh}\tau_k}
     {\sum_{h,k\in i}w_h d_{kh}},\qquad
e_i=
\frac{\sum_{h,k\in i}w_h d_{kh}\mathbf 1(\tau_k>0)}
     {\sum_{h,k\in i}w_h d_{kh}}.
\label{eq:sector-rates}
\end{equation}
Cette pondération conserve la composition effective de la consommation
nationale, mais $e_i$ reste un proxy de droit à déduction construit depuis la
demande finale des ménages. Il ne mesure ni le taux légal sur les intrants, déjà
porté par $\bar\tau_i$, ni la collecte lors de la vente finale, portée par
$\alpha_{kh}$. Pour une branche sans produit observé, le taux par défaut est
18~\%, sauf services structurellement exonérés (finance, immobilier,
administration, éducation, santé et services domestiques), auxquels sont
attribués un taux et une part taxable nuls.

Une sensibilité juridique remplace ce proxy par un indicateur binaire
$e_i^{L}$. Il vaut zéro pour les produits structurellement exonérés
\texttt{K38}, \texttt{L39}, \texttt{O42}, \texttt{P43}, \texttt{Q44},
\texttt{T47} et \texttt{U48}, et un pour les autres produits. Ce classement
simplifié ne prétend pas reproduire toutes les exceptions du CGI; il mesure ce
que devient la TVA incorporée lorsque le droit à déduction dépend d'une
classification juridique plutôt que de la composition des achats de l'enquête.

\subsection{TVA non déductible et propagation dans les chaînes de production}

La TVA sur intrants n'est un coût que pour la fraction de production qui ne
dispose pas du droit à déduction. La charge de premier tour du produit $i$ est
donc
\begin{equation}
b=(\mathbf 1-e)\odot(A^{d\prime}+A^{m\prime})\bar\tau,
\label{eq:first-round}
\end{equation}
où $\odot$ désigne le produit terme à terme. Une fois incorporée au coût d'un
produit domestique, cette charge est transmise à toutes les productions qui
l'utilisent en aval. Le vecteur complet de TVA enchâssée est la solution de
\begin{equation}
c=b+A^{d\prime}c
  =(I-A^{d\prime})^{-1}b.
\label{eq:leontief-vat}
\end{equation}
Les intrants importés contribuent au premier tour, mais seule la production
domestique est itérée dans la chaîne ivoirienne. Cette distinction correspond
au mécanisme économique : un produit final dans une branche devient la
consommation intermédiaire d'une autre, qui transmet à son tour le coût déjà
incorporé.

Pour ne pas compter deux fois la même dépense TTC, les composantes directe et
enchâssée utilisent un dénominateur commun. Pour une dépense $d$, un taux direct
$\tau$, un taux enchâssé $c$ et une part formelle $\alpha$, nous calculons
\begin{align}
T^d(d)&=\alpha d\frac{\tau}{1+\tau+c},\\
T^e(d)&=\alpha d\frac{c}{1+\tau+c}
       +(1-\alpha)d\frac{c}{1+c}.
\label{eq:io-incidence}
\end{align}
Le vecteur $c$ est défini ici comme une part additive du prix final hors marges,
au même niveau que la TVA finale : le prix observé est décomposé entre base,
$\tau$ et $c$. Il ne s'agit pas d'un multiplicateur de coût placé avant la TVA.
Cette convention explique l'absence de terme croisé $\tau c$ dans le
dénominateur $1+\tau+c$. Une convention multiplicative ferait apparaître ce
terme et constituerait un autre modèle de prix. Une transaction finale
informelle ne supporte pas la TVA directe, mais son prix peut contenir une taxe
incorporée plus en amont. La figure~\ref{fig:vat-chain} résume cette mécanique.

\begin{figure}[t]
\centering
\begin{tikzpicture}[
  node distance=12mm and 18mm,
  box/.style={draw, rounded corners=2pt, align=center, text width=3.05cm,
              minimum height=15mm, fill=gray!4, inner sep=3pt},
  arr/.style={-{Latex[length=2mm]}, thick}
]
\node[box] (input) {Intrant acheté\\TVA facturée au taux $\bar\tau$};
\node[box, right=of input] (producer) {Producteur sans droit à déduction\\TVA d'amont non récupérée};
\node[box, right=of producer] (embedded) {Produit issu de la branche\\charge de premier tour $b$};
\node[box, below=of embedded] (downstream) {Branche située en aval\\le produit devient un intrant};
\node[box, left=of downstream] (final) {Bien ou service final\\charge incorporée totale $c$};
\draw[arr] (input) -- (producer);
\draw[arr] (producer) -- (embedded);
\draw[arr] (embedded) -- (downstream);
\draw[arr] (downstream) -- (final);
\draw[arr, dashed] (downstream.east) .. controls +(1.1,0) and +(1.1,0) ..
  (embedded.east);
\end{tikzpicture}
\caption{Propagation de la TVA non déductible dans une chaîne de production.
Le produit d'une branche peut devenir l'intrant d'une autre; l'inverse de
Leontief additionne tous les tours domestiques.}
\label{fig:vat-chain}
\end{figure}

\subsection{Accises et droits de douane}

Les accises et droits de douane sont retirés de la dépense TTC en respectant
l'ordre légal de la cascade. Chaque poste de consommation est relié à un produit
du TRE, à une bande du TEC et, s'il est accisable, à un taux 2021. La part
importée du produit $k$ est
$m_k=M_k/(M_k+Q_k)$, à partir des importations et de la production domestique du
TRE; la part de marges commerciales et de transport est notée $\mu_k$. Pour une
dépense $d_{kh}$, la base hors marges est $B_{kh}=d_{kh}(1-\mu_k)$. Si
$\epsilon_k$ est le taux d'accise, $\delta_k$ la bande TEC et $r_k$ le taux de
TVA, les deux charges directement contenues dans l'achat sont
\begin{align}
A_{kh}&=B_{kh}\frac{\epsilon_k}{(1+\epsilon_k)(1+r_k)},\\
T^{\mathrm{douane}}_{kh}&=B_{kh}\frac{\delta_km_k}
 {(1+\delta_km_k)(1+\epsilon_k)(1+r_k)}.
\label{eq:other-indirect}
\end{align}
Le facteur $(1+r_k)$ suppose que la couche de TVA payée en amont reste contenue
dans le prix hors marges, même lorsque la vente finale n'est pas soumise à la
TVA directe. Le coefficient $\alpha_{kh}$ ne modifie donc pas cette couche : il
porte seulement sur la collecte à la vente finale. Si l'informalité supprimait
aussi la TVA d'amont du prix observé, ces formules sous-estimeraient les accises
et droits de douane extraits des achats concernés. Cette décomposition évite de
retirer séparément chaque taxe de son propre montant TTC : le calcul chiffre à
28,6 milliards de FCFA le chevauchement de TVA qui résulterait de cette
opération naïve.

Le scénario central pose une transmission unitaire. Les droits de porte sont
acquittés à l'importation et les accises à la production ou à la sortie de
régime suspensif; l'informalité du détaillant ne les annule donc pas. Pour le
tabac et l'alcool seulement, les scénarios S2 et S3 appliquent les coefficients
d'achat formel comme sensibilités de contrebande. Les carburants restent à
transmission unitaire. Les paramètres et références budgétaires sont documentés
avec leur source et sont rapprochés des encaissements de la loi de règlement
2021 \citep{civreglement2021}. Le calcul alloue les taxes directement
contenues dans les achats des ménages; leur propagation amont par Leontief reste
une extension.

\subsection{Impôts directs, assujettissement fiscal et CMU}

Nous définissons un scénario central d'assujettissement multi-critères. Son
objectif n'est pas de classer tout l'emploi en formel ou informel, mais
d'identifier les rémunérations salariales pour lesquelles l'enquête observe au
moins un indice de relation d'emploi traçable et donc de retenue fiscale
plausible. L'emploi principal est retenu lorsqu'une rémunération positive et au
moins un mois de travail sont déclarés, conjointement avec au moins un des
éléments suivants : droit à des congés payés, nombre positif de jours de congé
annuel pris, congé maladie, cotisation à la CGRAE ou à la CNPS, congé de
maternité ou de paternité, ou bulletin de salaire.

Pour l'emploi secondaire, le questionnaire ne recueille pas le même ensemble
d'avantages contractuels. Nous retenons donc les rémunérations positives des
emplois appartenant aux grands groupes de direction, professions
intellectuelles et scientifiques, professions intermédiaires, employés de
bureau et forces armées, ainsi qu'aux catégories professionnelles détaillées
associées à un emploi salarié. Le salaire assujetti est la somme annualisée des rémunérations principale et
secondaire retenues. Cette assiette salariale est construite une seule fois et
reste identique dans toutes les comparaisons. Les scénarios diffèrent seulement
par le critère d'assujettissement : le scénario strict exige simultanément un
bulletin de salaire et une cotisation sociale; le scénario élargi exige l'un ou
l'autre. Ils sont donc des sous-ensembles emboîtés dans le scénario central, et
une différence de masse fiscale ne peut plus provenir d'une autre définition du
salaire.

Les non-réponses au bulletin et à la cotisation sont conservées avant tout
recodage. Elles ne sont jamais considérées comme une preuve d'emploi formel. La
règle centrale peut néanmoins retenir un salarié lorsqu'un autre indice positif
est observé, par exemple un congé payé ou une cotisation déclarée. Dans les
critères strict et élargi, un marqueur manquant ne satisfait pas le critère. Le
diagnostic de non-réponse publie les effectifs concernés mais conserve
exactement le classement central; il ne crée aucun assujetti supplémentaire.

Cette définition est une règle d'imputation fiscale propre à l'étude, pas une
mesure générale de l'emploi formel au sens de l'OIT. À titre diagnostique
seulement, nous utilisons le critère observable le plus proche de la
21\textsuperscript{e} Conférence internationale des statisticiens du travail
(CIST) :
un emploi salarié principal accompagné d'une cotisation déclarée à la CGRAE ou
à la CNPS. Une variante ajoute le bénéfice conjoint de congés payés et maladie
\citep{ilo2023}. Ce diagnostic est exporté séparément et ne modifie pas le
scénario fiscal central.

Les barèmes CGI sont appliqués mensuellement. Si $W$ est le salaire brut
imposable, l'IS vaut $0,012W$; la CN est la somme des prélèvements marginaux par
tranche sur $0,8W$. La base soumise à l'IGR est notée
$B_{\mathrm{IGR}}=0,8(W-IS-CN)$ et le quotient familial vaut
$Q=B_{\mathrm{IGR}}/N$. L'impôt par part du barème à dix tranches est multiplié par $N$, puis
réduit de 15~\%. Le nombre de parts dépend de l'état civil, augmente de 0,5 par
enfant de moins de 18 ans, est plafonné à cinq, et les enfants sont affectés au
salaire assujetti le mieux rémunéré du ménage. Les égalités de salaire sont
départagées par l'ordre stable de l'identifiant individuel. Une sensibilité
partage également les demi-parts d'enfants entre les assujettis du ménage. La
CNPS salariale vaut 6,3~\% du
salaire dans la limite mensuelle de 1~647~315 FCFA.

Pour la CMU, le vivier observable comprend les personnes qui déclarent une
assurance maladie financée par un organisme public. Un tirage reproductible
(graine 74375322) ordonne séparément les pauvres
et les non-pauvres; les cibles sont 75~000 participants pauvres non contributifs
et 1~425~000 participants non pauvres contributifs. La contribution vaut 1~000
FCFA par mois pour les adultes du second groupe. Dans la livraison EHCVM
utilisée, le vivier observable ne permet d'atteindre que 50~641 participants
pauvres et 885~822 non-pauvres, dont 537~231 adultes contributeurs; tous les
admissibles sont alors retenus et l'écart aux cibles est publié. Une variante
statutaire applique 1~000 FCFA mensuels à tous les adultes afin de borner la
masse potentielle.

\subsection{Pensions, transferts directs et filets sociaux}

Le module recense ou impute les sommes monétaires que les ménages reçoivent des
administrations publiques. Combinées aux impôts directs et aux cotisations,
elles permettent de remonter du revenu disponible observé vers les ressources
dont le ménage disposait avant l'intervention publique. Le module distingue
quatre catégories principales de paiements publics. Premièrement, le PSSN est
reconstitué par un test indirect de niveau de vie, ou
PMT. Le registre administratif n'étant pas apparié à l'enquête, cette procédure
ne cherche pas à reconnaître les bénéficiaires réels un par un; elle reproduit
une règle de ciblage transparente avec les informations disponibles dans
l'EHCVM. Cette logique générale --- prédire le niveau de vie à partir de
caractéristiques observables --- suit la formulation classique des tests
indirects de niveau de vie \citep{groshbaker1995}. Le projet ivoirien prévoit
lui-même un dispositif de ciblage et un registre social
\citep{worldbankpssn2015}; la formule reconstruite ici ne doit toutefois pas être
confondue avec sa formule administrative non observée. Une régression linéaire
pondérée par le poids du ménage explique le
logarithme du revenu disponible par personne. Les variables explicatives sont
le logarithme de la taille du ménage; le milieu urbain ou rural; la région;
l'âge du chef et son carré; le sexe, le niveau d'éducation et le statut
matrimonial du chef; puis les nombres de membres âgés de 0 à 5 ans, de 6 à
10 ans, de 11 à 15 ans et de 65 ans ou plus, ainsi que le nombre de personnes
ayant un handicap déclaré. Le milieu, la région et la taille proviennent du
fichier ménage; les caractéristiques du chef et la composition démographique,
du fichier individus. Lorsque l'âge du chef manque, il est remplacé par l'âge
médian observé; une caractéristique catégorielle manquante prend la modalité de
référence et un effectif démographique manquant vaut zéro.

Le niveau de vie prédit est obtenu en exponentiant la prédiction du modèle. Les
ménages sont ensuite classés de la prédiction la plus faible à la plus élevée;
l'identifiant ménage départage les égalités, ce qui rend la sélection
reproductible. Le cumul des poids ménage fixe deux coupures : une première au
plus près de 177~000 ménages pour la cohorte présente toute l'année, puis une
seconde au plus près de 192~000 pour inclure la cohorte partielle. Comme un
ménage enquêté représente un nombre entier de ménages pondérés et n'est pas
fractionné à la frontière, le calcul retient 177~142 ménages en année complète
et 14~625 en cohorte partielle, soit 191~767 au total et un écart de 233 à la
cible. Cette différence de couverture est publiée; aucun poids n'est modifié
pour la faire disparaître.

L'allocation statutaire vaut 36~000 FCFA par trimestre. Quatre versements sont
attribués à la première cohorte et trois à la seconde. La masse obtenue avant
ajustement est multipliée par 1,000739 afin de retrouver exactement
27,108 milliards de FCFA, montant exécuté documenté pour 2021
\citep{civperformance2021}. Le calage porte donc sur la masse monétaire, pas sur
le rang PMT ni sur la distribution entre déciles. Dans le volet de l'enquête consacré aux programmes sociaux, 24~989 ménages
pondérés déclarent recevoir le programme monétaire, sans en indiquer le montant;
seuls 972 appartiennent aussi à la sélection PMT. La collecte ne permet pas de
vérifier le programme administratif exact ni le montant reçu; ces déclarations
ne remplacent donc pas le score dans le scénario central.

Une robustesse « déclarés d'abord » retient en priorité tous ces ménages, puis
complète la cible par ordre croissant du score PMT. Elle conserve les deux
cohortes, ne fractionne aucun poids à la coupure et cale de nouveau la masse sur
27,108 milliards de FCFA. Elle retient 192~544 ménages pondérés, conserve les
24~989 déclarants et attribue 79,5~\% des bénéficiaires à des ménages sous le
seuil, contre 85,9~\% dans le PMT central. L'annexe statistique publie les
coefficients, les variables, les coupures, les facteurs de calage et les profils
par décile des deux allocations.

La mesure des paiements publics n'attribue aucun montant monétaire aux aides
alimentaires ni aux
prestations de santé, car elles sont fournies en nature. Elle n'ajoute pas non
plus les travaux à haute intensité de main-d'œuvre (HIMO) aux paiements publics :
ils rémunèrent une activité et peuvent déjà apparaître dans les revenus du
travail déclarés. Les participations observées à ces programmes restent
documentées pour les modules qui peuvent les valoriser.

Deuxièmement, les bourses correspondent aux montants directement déclarés par
les élèves du secondaire général public et du supérieur public.
Troisièmement, les prestations familiales appliquent les barèmes CNPS aux
enfants et naissances observés des ménages comprenant un travailleur couvert.
Quatrièmement, la dépense annuelle d'accidents du travail est distribuée comme
valeur attendue entre travailleurs présumés couverts. Cette couverture n'utilise
pas la déclaration d'assurance maladie : elle est approchée par la cotisation
retraite déclarée et la durée de l'emploi. Les pensions de retraite, de
survivant et d'invalidité sont directement observées parmi les revenus
déclarés. Les pensions alimentaires et les sommes reçues de parents ou d'autres ménages
sont présentées séparément dans les diagnostics. Elles ne sont pas comptées
comme des paiements publics, puisqu'il s'agit d'argent transféré entre agents
privés et non de ressources versées par une administration.

Dans le scénario PDI central, les pensions contributives sont déjà comprises
dans les ressources acquises avant l'intervention publique et les cotisations
retraite sont considérées comme une épargne obligatoire. Pour remonter du revenu
disponible vers ces ressources initiales, le calcul réintègre donc seulement les
prélèvements non liés à la retraite qui sont mesurés ici, notamment l'IRPP et la
CMU, puis retire les paiements monétaires des programmes publics. Dans la
robustesse PGT, le même calcul compte au contraire les pensions parmi les
paiements publics reçus et les cotisations retraite parmi les prélèvements
supportés. Les deux reconstructions retrouvent le même revenu disponible à
moins de $10^{-6}$ FCFA. Cette double présentation suit le manuel CEQ et évite
qu'une convention comptable sur les pensions détermine silencieusement le
diagnostic distributif \citep{lustig2022}.

\subsection{Réductions publiques de prix sur l'électricité, l'eau et les carburants}

Le calcul rapproche les dépenses détaillées d'électricité, d'eau et de carburant
des informations sur le logement et l'abonnement aux réseaux. Pour l'eau, il
distingue l'achat facturé directement par le réseau de l'achat auprès d'un
revendeur. Le raccordement, l'inclusion de la facture dans le loyer et le type
de compteur servent de contrôles de cohérence, mais un ménage sans compteur
individuel n'est pas automatiquement écarté.

Pour l'électricité, le rapport de l'ANARE-CI documente une aide d'exploitation
de 8,69 milliards de FCFA, environ 30~\% d'abonnés au tarif social, un tarif
moyen hors taxes de 51 FCFA par kWh pour ces abonnés et de 70 FCFA pour le tarif
domestique général \citep{anare2021}. Parmi les ménages ayant une facture
positive, les 30~\% aux factures pondérées les plus faibles sont classés dans
le groupe social. La facture est convertie en quantité avec le tarif
correspondant, puis les 8,69 milliards sont répartis au prorata de cette
quantité. Le montant national est ainsi retrouvé exactement; des variantes à
20 et 40~\% d'abonnés sociaux testent la sensibilité de la frontière.

Pour l'eau, le barème trimestriel est inversé tranche par tranche : 235 FCFA par
mètre cube pour les 18 premiers mètres cubes, 367,3 jusqu'à 90, puis 586,8
jusqu'à 300. La réduction de prix centrale est l'écart entre le tarif domestique
et le tarif social appliqué aux 72 premiers mètres cubes annuels. Les achats
auprès d'un revendeur ne reçoivent rien dans le scénario central, faute de tarif
observé; une marge de 20~\% de la dépense forme la borne haute. Enfin, le prix à
la pompe 2021 permet d'estimer les litres achetés, mais aucune série annuelle de
prix de parité suffisamment auditée ne permet d'identifier un soutien au
carburant. Son montant central est donc nul. Une borne haute de 80 FCFA par
litre, fondée sur la structure de prix postérieure et explicitement séparée de
2021, n'est utilisée qu'en robustesse \citep{bceao2023}. Cette prudence évite de
transformer une hypothèse de prix en dépense publique observée. La méthode suit
le traitement CEQ des subventions à la consommation \citep{lustig2022,younger2017}.

\subsection{Services publics d'éducation}

L'analyse d'incidence des dépenses d'éducation suit l'approche du coût moyen par
utilisateur \citep{demery2003,lustig2022,akim2020}. Les informations scolaires
de l'EHCVM permettent d'observer la fréquentation pendant l'année, le niveau
d'études, le caractère public ou privé de l'établissement et les paiements
effectués par les familles. Les niveaux maternel et primaire sont réunis; le
secondaire général, le secondaire technique, le postsecondaire professionnel et
le supérieur sont traités séparément.

Les dépenses exécutées 2021 proviennent du Rapport général sur la performance :
587,505 milliards pour le préscolaire et le primaire, 378,934 pour le secondaire
général, 13,052 pour le secondaire technique, 91,771 pour la formation
professionnelle et 283,609 pour l'enseignement supérieur et les œuvres
universitaires \citep{civperformance2021}. Dans chaque niveau, le budget est
divisé par le nombre pondéré d'élèves du public observés dans l'enquête. Chaque
élève reçoit ce coût public moyen; la somme pondérée retrouve donc exactement le
budget du niveau. Le bénéfice net retranche les frais d'inscription et les
contributions scolaires directement payés. Deux robustesses utilisent les seules
dépenses courantes, puis retranchent l'ensemble des dépenses scolaires déclarées.
Cette valorisation mesure un service financé par l'État, pas la qualité de
l'enseignement ni son rendement futur.

\subsection{Services publics de santé}

Le calcul de santé écarte les déclarations relatives à l'assurance maladie :
leur couverture est trop incomplète pour déterminer le bénéfice central. Les
informations retenues décrivent les maladies et les consultations des trente
derniers jours, les dépenses de soins des trois derniers mois, ainsi que les
hospitalisations des douze derniers mois, leur nombre, le type de prestataire et
les frais acquittés. Les structures publiques sont identifiées à partir de la
catégorie de prestataire déclarée.

Le rappel de trente jours rend l'utilisation ambulatoire observée volatile. Le
scénario central attribue donc à chaque personne la dépense publique moyenne
attendue pour des personnes comparables. La probabilité pondérée d'une
consultation publique est calculée dans des cellules croisant six classes d'âge
(0--4, 5--14, 15--24, 25--44, 45--64 et 65 ans ou plus), le sexe et le milieu.
Lorsqu'une cellule contient moins de 30 observations, le taux âge--sexe national
la remplace. Ce taux est annualisé par douze. Il ne repose ni sur une adhésion à
la CMU ni sur une déclaration d'assurance. Les hospitalisations sont, elles,
attribuées selon les séjours publics effectivement déclarés sur douze mois; un
séjour représente dix unités de ressources ambulatoires dans la spécification centrale, contre cinq
et vingt en robustesse.

Le budget central de 108,527 milliards additionne les actions exécutées de
médicaments et intrants et de prise en charge médicale. Il est divisé par le
nombre pondéré d'unités ambulatoires attendues et d'unités hospitalières, puis
entièrement réparti. La borne budgétaire large de 141,737 milliards ajoute le
reste du programme santé \citep{civperformance2021}. Le bénéfice net retranche
les dépenses de soins annualisées par quatre et les frais hospitaliers payés
aux structures publiques. L'allocation intégrale aux consultations observées,
le budget large et les poids hospitaliers 5 et 20 sont publiés en robustesse.
Il s'agit d'une incidence budgétaire au coût de production, non d'une mesure de
l'état de santé, de la qualité des soins ou d'une valeur dérivée de l'assurance
\citep{demery2003,lustig2022}.

\subsection{Assemblage des revenus, contribution distributive et appauvrissement fiscal}

Le revenu consommable retranche du revenu disponible la TVA directe et la TVA
non déductible incorporée dans les prix, les accises et les droits de douane,
puis ajoute les réductions publiques de prix. Le revenu final ajoute ensuite
les bénéfices nets d'éducation et de santé. Les mêmes opérations sont réalisées
sous PDI et PGT; les deux conventions diffèrent avant le revenu disponible mais
retrouvent le même revenu disponible observé. Les identités sont contrôlées
ménage par ménage à une tolérance de $10^{-6}$ FCFA. Les résultats centraux
utilisent S3 pour la TVA, la transmission complète des accises et droits de
douane, les réductions de prix centrales et les bénéfices en nature nets des
paiements directs.

Pour expliquer la variation du Gini sans dépendre d'un ordre arbitraire, les
instruments sont réunis en six groupes : prélèvements directs, paiements publics
directs, impôts indirects, réductions de prix, éducation et santé. La valeur de
Shapley moyenne la contribution marginale de chaque groupe sur les 720 ordres
possibles, calculés exactement à partir des 64 sous-ensembles
\citep{shorrocks2013}. Cinquante réplications Rao--Wu recalculent la
décomposition complète; les autres indicateurs utilisent 500 réplications.

L'appauvrissement fiscal compare le revenu primaire monétaire au revenu
consommable. Les services d'éducation et de santé ne sont pas convertis en
argent pour déterminer un franchissement du seuil. Pour un seuil $z$, la perte
sous le seuil est
$\max\{\min(Y_P,z)-\min(Y_C,z),0\}$ et le gain est l'expression symétrique. Leur
différence pondérée, divisée par $z$, retrouve exactement la variation de la
profondeur de pauvreté. La part appauvrie, les nouveaux pauvres, les gains sous
le seuil et leurs montants sont rapportés séparément, conformément à
\citet{higgins2016}; une grille allant de 0,5 à 2 fois le seuil national teste la
dominance du résultat.

\subsection{Pauvreté, progressivité et inférence de sondage}

Les indices de \citet{fgt1984} sont
\begin{equation}
P_a=\frac{\sum_h w_hn_h
\left[\max\{(z_h-y_h)/z_h,0\}\right]^a}
{\sum_h w_hn_h},\qquad a\in\{0,1,2\},
\label{eq:fgt}
\end{equation}
où $y_h$ est le concept de revenu par tête, $z_h$ le seuil spécifique, $w_h$ le
poids ménage et $n_h$ sa taille. Le Kakwani est
$K=C_T-G_Y$, différence entre la concentration de l'impôt et le Gini du revenu
de classement \citep{kakwani1977}. Le Reynolds--Smolensky est rapporté selon la
convention $RS=G_{\text{après}}-G_{\text{avant}}$ : une valeur négative
indique une réduction des inégalités. Les rangs fractionnaires pondérés suivent
la règle du point médian.

L'incertitude respecte le plan stratifié en grappes. Dans chaque strate $h$
contenant $m_h$ unités primaires, le bootstrap Rao--Wu tire $m_h-1$ grappes avec
remise et multiplie leurs poids par $m_h/(m_h-1)$ \citep{raowu1988}. Les
intervalles percentiles reposent sur 500 réplications. Les mêmes poids sont
utilisés pour comparer TRE courant, TRE constant et ICIO, de sorte que les
intervalles sur leurs différences incorporent la covariance entre scénarios.

\subsection{Taxation de la vente finale et recyclage des recettes}

Pour les produits initialement à taux nul et inclus dans l'ensemble réformé
$\mathcal K$, la dépense observée est une base hors TVA. La charge additionnelle
nominale est donc
\begin{equation}
\Delta T_h=0,09\sum_{k\in\mathcal K}\alpha_{kh}d_{kh}.
\label{eq:reform-tax}
\end{equation}
La réforme porte sur la TVA facturée à la vente finale. Les coefficients
techniques et le vecteur $c$ de taxe déjà incorporée restent constants : nous
n'attribuons donc pas mécaniquement le nouveau taux de 9~\% à la composante
enchâssée. Une modification des droits à déduction sectoriels exigerait un
nouveau contrefactuel I/O et constitue une réforme distincte.
Nous simulons les ensembles agricole et large sous transmission intégrale, S2 et S3. Le revenu
après réforme brute est
$Y'_{Dh}=Y_{Dh}-\mathrm{def}_h\Delta T_h/n_h$. Pour le scénario à budget constant, la
recette $\mathcal R=\sum_hw_h\Delta T_h$ est divisée par la population
$N=\sum_hw_hn_h$. Chaque personne reçoit $b=\mathcal R/N$, si bien que
\begin{equation}
Y''_{Dh}=Y_{Dh}-\mathrm{def}_h\frac{\Delta T_h}{n_h}
          +\mathrm{def}_h b,\qquad
\sum_hw_hn_hb=\mathcal R.
\label{eq:recycling}
\end{equation}
La recette et $b$ sont recalculés à chaque réplication Rao--Wu. Le résultat
compare ainsi une réforme sans emploi des fonds et un paquet fiscal
budgétairement neutre; il ne suppose ni coût administratif ni ciblage parfait.

%=============================================================================
\section{Résultats}
\label{sec:resultats}

\subsection{TVA directe : une charge croissante avec le niveau de vie}

Le tableau~\ref{tab:vat-scenarios} rassemble les trois hypothèses de taxation
effective. Dans le scénario central S3, la TVA directe simulée atteint 693
milliards de FCFA. Le taux effectif moyen passe de 2,01~\% du revenu disponible
au premier décile à 6,13~\% au dixième. Cette pente combine la composition des
paniers et l'augmentation de la probabilité d'achat formel avec le niveau de
vie. Le Kakwani vaut $0,165$ et son intervalle à 95~\%, $[0,157;0,173]$, exclut
nettement zéro.

\begin{table}[t]
\centering
\begin{threeparttable}
\caption{TVA directe selon l'hypothèse de taxation effective}
\label{tab:vat-scenarios}
\begin{tabular}{lrrrr}
\toprule
Hypothèse & Masse & Kakwani & IC à 95~\% & $RS$ \\
           & (mds FCFA) & & & \\
\midrule
Strict : transmission intégrale & 1~387 & 0,076 & [0,067 ; 0,084] & $-0,0069$ \\
S2 : COICOP $\times$ milieu     &   739 & 0,132 & [0,123 ; 0,140] & $-0,0064$ \\
S3 : COICOP $\times$ décile     &   693 & 0,165 & [0,157 ; 0,173] & $-0,0076$ \\
\bottomrule
\end{tabular}
\begin{tablenotes}\footnotesize
\item \emph{Notes} : masses issues du TRE courant afin de conserver la même
décomposition TTC que le modèle I/O. Kakwani et Reynolds--Smolensky utilisent
$Y_D$ comme revenu de classement. Intervalles Rao--Wu, 500 réplications.
Source : calculs des auteurs à partir de l'EHCVM 2021.
\end{tablenotes}
\end{threeparttable}
\end{table}

La transmission intégrale double approximativement la masse par rapport à S3,
mais ne change pas le signe de la progressivité : même sans informalité
modélisée, les produits taxés occupent une part plus forte des paniers du haut de
la distribution. En revanche, le paramétrage de l'informalité accentue fortement
la pente : sous transmission intégrale, le taux effectif va de 6,48 à 10,4~\% entre les déciles
extrêmes; sous S3, il va de 2,01 à 6,13~\%. La figure~\ref{fig:vat-scenarios}
montre que l'écart entre S2 et S3 est surtout concentré au bas de la
distribution, tandis que les deux profils convergent au sommet.

\begin{figure}[t]
\centering
\includegraphics[width=0.88\textwidth]{fig1_eff_vat_three_scenarios.png}
\caption{Taux effectif de TVA directe par décile et scénario de taxation
effective. S3 est le scénario central; S2 et la transmission intégrale sont des sensibilités.
Source : calculs des auteurs à partir de l'EHCVM 2021.}
\label{fig:vat-scenarios}
\end{figure}

\subsection{Une taxe progressive peut néanmoins accroître la pauvreté}

La progressivité est une propriété relative de la charge le long de la
distribution; elle ne garantit pas que le niveau de vie des pauvres augmente.
Dans le scénario S3, la TVA directe fait passer $P_0$ de 37,47 à 39,81~\%, soit
$+2,33$ points avec un intervalle $[1,88;2,85]$. L'ajout de la TVA enchâssée
porte $P_0$ à 40,96~\%, soit $+3,48$ points $[2,96;4,03]$. Le même mécanisme
apparaît sous S2 et sous transmission intégrale, avec une amplitude croissante
de la charge.

\begin{table}[t]
\centering
\begin{threeparttable}
\caption{Pauvreté avant TVA et après TVA directe et totale}
\label{tab:fgt-io}
\begin{tabular}{llrrr}
\toprule
Scénario & Concept & $P_0$ (\%) & $\Delta P_0$ (pt) & IC à 95~\% \\
\midrule
Référence & $Y_D$, avant TVA & 37,47 & 0,00 & -- \\
S3 & TVA directe & 39,81 & 2,33 & [1,88 ; 2,85] \\
S3 & TVA directe + enchâssée & 40,96 & 3,48 & [2,96 ; 4,03] \\
S2 & TVA directe + enchâssée & 41,42 & 3,94 & [3,37 ; 4,57] \\
Strict & TVA directe + enchâssée & 44,49 & 7,01 & [6,32 ; 7,73] \\
\bottomrule
\end{tabular}
\begin{tablenotes}\footnotesize
\item \emph{Notes} : TRE 2023 courant. Les poids de population multiplient le
poids du ménage par le nombre de ses membres. Le seuil national de pauvreté de
l'EHCVM 2021 est de 369~516 FCFA par personne et par an.
Intervalles Rao--Wu, 500 réplications conjointes. Source : calculs des auteurs à partir de l'EHCVM 2021.
\end{tablenotes}
\end{threeparttable}
\end{table}

Dans la spécification centrale, 191~467 ménages, représentant 1~038~623
personnes, se situent au-dessus du seuil dans $Y_D$ et en dessous dans le revenu
après TVA totale; 47,8~\% de ces nouveaux pauvres vivent en milieu rural. Ces
basculements ne signifient pas qu'un million de personnes supportent toute la
charge : ils indiquent que la marge de revenu de ces personnes au-dessus du
seuil était inférieure à la taxe simulée.

\subsection{La TVA enchâssée et la robustesse au choix de matrice}

Le TRE courant produit, sous S3, 214 milliards de FCFA de TVA enchâssée, contre
693 milliards de TVA directe. La composante amont représente donc 23,6~\% de la
charge totale. Elle ne se confond ni avec la ligne de TVA non déductible du TRE,
qui atteint 1~611 milliards pour l'ensemble des emplois de l'économie, ni avec
une nouvelle taxe facturée au consommateur : il s'agit de la part de coût
incorporée à la consommation des ménages après propagation et concordance.

\begin{table}[t]
\centering
\begin{threeparttable}
\caption{Sensibilité de la TVA totale à la matrice et au droit à déduction, scénario S3}
\label{tab:io-robustness}
\begin{tabular}{lrrrr}
\toprule
Spécification & TVA directe & TVA enchâssée & TVA totale & $\Delta P_0$ \\
        & \multicolumn{3}{c}{(milliards de FCFA)} & (point) \\
\midrule
TRE 2023, prix courants  & 693 & 214 & 907 & 3,48 \\
TRE courant, $e_i^L$ juridique & 697 & 22 & 719 & -- \\
TRE 2023, prix constants & 693 & 220 & 913 & 3,56 \\
OCDE ICIO 2020           & 691 & 212 & 903 & 3,56 \\
\bottomrule
\end{tabular}
\begin{tablenotes}\footnotesize
\item \emph{Notes} : le TRE courant est la référence. La ligne juridique pose
un droit à déduction complet hors sept produits structurellement exonérés; son
effet sur la pauvreté n'est pas recalculé. Le tableau du TRE à prix constants
porte le millésime 2023 mais ne documente pas l'année de base de ses prix; il
est donc utilisé comme test de robustesse et non pour interpréter une évolution
en volume. L'écart de
$\Delta P_0$ entre TRE courant et ICIO est $-0,072$ point, IC à 95~\%
$[-0,197;0,031]$. La couverture pondérée de la concordance est 92,5~\% pour le TRE et
95,2~\% pour ICIO. Source : calculs des auteurs à partir de l'EHCVM 2021.
\end{tablenotes}
\end{threeparttable}
\end{table}

Les trois estimations fondées sur des matrices différentes sont proches. Le TRE
constant augmente la TVA enchâssée de 5,8 milliards par rapport au TRE courant.
ICIO donne une masse totale inférieure de 4,1 milliards, mais un effet pauvreté
légèrement supérieur en raison de la distribution sectorielle de la charge.
L'intervalle sur la différence TRE--ICIO contient zéro.

La sensibilité au droit à déduction est beaucoup plus forte. Le classement
juridique binaire réduit la TVA enchâssée de 214,4 à 22,3 milliards et la charge
totale de 907,3 à 719,0 milliards; la TVA directe passe de 692,9 à 696,7
milliards parce que le dénominateur commun contient moins de TVA incorporée. Ce
contrôle est volontairement grossier, mais il montre que le résultat central ne
peut pas être présenté comme robuste au choix de $e_i$. Une nomenclature
administrative des droits à déduction par produit est nécessaire pour resserrer
cet intervalle structurel. La figure~\ref{fig:io} présente les niveaux de
pauvreté correspondant aux seules variantes de matrice.

\begin{figure}[t]
\centering
\includegraphics[width=0.88\textwidth]{fig_poverty_io_impact.png}
\caption{Pauvreté avant TVA, après TVA directe et après TVA totale. Le panneau
central repose sur le TRE 2023 courant et S3; les autres barres encadrent les
hypothèses de transmission et de matrice. Source : calculs des auteurs à partir de l'EHCVM 2021.}
\label{fig:io}
\end{figure}

\subsection{Accises et droits de douane : une charge additionnelle diffuse}
\label{sec:res-other-indirect}

Les achats des ménages contiennent 67,9 milliards de FCFA d'accises et 185,1
milliards de droits de douane. La décomposition des accises donne 36,5 milliards
sur les carburants, 13,6 sur l'alcool, 7,6 sur les cosmétiques, 7,1 sur le tabac
et 3,2 sur les boissons non alcoolisées. Le tableau~\ref{tab:other-indirect}
montre que les taux effectifs augmentent avec le niveau de vie, mais moins
fortement que pour la TVA S3.

\begin{table}[t]
\centering
\begin{threeparttable}
\caption{Accises et droits de douane directement contenus dans les achats}
\label{tab:other-indirect}
\begin{tabular}{lrrrr}
\toprule
Instrument & Masse & Taux D1 & Taux D10 & Kakwani \\
           & (mds FCFA) & \multicolumn{2}{c}{(\% de $Y_D$)} & \\
\midrule
Accises              &  67,9 & 0,292 & 0,548 & 0,102 \\
Droits de douane     & 185,1 & 0,946 & 1,348 & 0,059 \\
Ensemble             & 253,0 & 1,238 & 1,896 & 0,070 \\
\bottomrule
\end{tabular}
\begin{tablenotes}\footnotesize
\item \emph{Notes} : IC à 95~\% Rao--Wu : $[0,071;0,130]$ pour les accises,
$[0,034;0,088]$ pour les droits de douane et $[0,049;0,096]$ pour l'ensemble.
D1 et D10 désignent les premier et dixième déciles. Les taux sont rapportés au
revenu disponible réel. Source : calculs des auteurs à partir de l'EHCVM 2021.
\end{tablenotes}
\end{threeparttable}
\end{table}

Les montants imputés aux ménages représentent 23,2~\% des encaissements administratifs
comparables d'accises et 33,3~\% des droits de douane. Cet écart est attendu :
les références, 292,9 et 555,1 milliards, couvrent aussi les achats des
entreprises et administrations, tandis que la simulation ne porte que sur la
consommation directement identifiable des ménages. La concordance TEC couvre
100~\% de la dépense pondérée, avec des zéros explicites pour les services; la
concordance TRE qui fournit marges et parts importées en couvre 92,3~\%.

Ajoutées isolément à la TVA directe et enchâssée, les deux taxes porteraient la
pauvreté de 37,47 à 44,72~\%, soit $+7,25$ points. Cette transition partielle
n'inclut ni les paiements publics directs ni les réductions de prix. La séquence
monétaire complète donne 42,01~\% au revenu consommable. La propagation amont
des accises et droits de douane reste absente dans les deux calculs. La
figure~\ref{fig:other-indirect} présente les profils de
charge par décile.

\begin{figure}[t]
\centering
\includegraphics[width=0.88\textwidth]{fig17_indirect_other_decile.png}
\caption{Taux effectifs des accises et droits de douane par décile. Les barres
reposent sur les charges directement contenues dans les achats des ménages.
Source : calculs des auteurs à partir de l'EHCVM 2021.}
\label{fig:other-indirect}
\end{figure}

\subsection{Prélèvements directs : forte progressivité, assiette limitée}
\label{sec:res-direct}

Le scénario central multi-critères identifie 1,087 million d'assujettis pondérés,
soit 42,7~\% des salariés rémunérés observés. Au moins un assujetti est présent
dans 14,6~\% des ménages. Le taux de couverture augmente de 3,46~\% des ménages
du premier décile à 29,1~\% dans le dernier. Cette concentration explique
l'essentiel du profil de l'impôt et rejoint le mécanisme de capacité fiscale
mis en évidence par \citet{jensen2022} : l'extension du salariat couvert par un
tiers déclarant élargit progressivement l'assiette de l'impôt sur le revenu.

\begin{table}[t]
\centering
\begin{threeparttable}
\caption{Impôts directs et cotisations, scénario central multi-critères}
\label{tab:direct}
\begin{tabular}{lrrrr}
\toprule
Décile & Ménages assujettis & Taux IRPP & Taux cotisations & Part IRPP \\
       & (\%) & (\% de $Y_D$) & (\% de $Y_D$) & (\% de la masse) \\
\midrule
1  & 3,46 & 0,036 & 0,134 & 0,09 \\
5  & 7,05 & 0,290 & 0,356 & 1,58 \\
10 & 29,1 & 3,10  & 1,84  & 66,6 \\
\midrule
\multicolumn{2}{l}{IRPP} & \multicolumn{3}{r}{225,5 milliards FCFA} \\
\multicolumn{2}{l}{CNPS} & \multicolumn{3}{r}{155,9 milliards FCFA} \\
\multicolumn{2}{l}{CMU contributive observée} & \multicolumn{3}{r}{6,45 milliards FCFA} \\
\multicolumn{2}{l}{Kakwani IRPP / cotisations} & \multicolumn{3}{r}{0,452 / 0,344} \\
\bottomrule
\end{tabular}
\begin{tablenotes}\footnotesize
\item \emph{Notes} : l'IC à 95~\% du Kakwani est $[0,410;0,488]$ pour l'IRPP
et $[0,302;0,380]$ pour les cotisations. Les variantes d'assujettissement sont
détaillées dans le tableau~\ref{tab:direct-scenarios}. Source : calculs des
auteurs à partir de l'EHCVM 2021.
\end{tablenotes}
\end{threeparttable}
\end{table}


\begin{table}[t]
\centering
\begin{threeparttable}
\caption{Assiette et prélèvements selon le critère d'assujettissement}
\label{tab:direct-scenarios}
\begin{tabular}{lrrrr}
\toprule
Scénario & Assujettis & Masse salariale & IRPP & Cotisations \\
 & (pondérés) & \multicolumn{3}{c}{(milliards de FCFA)} \\
\midrule
Central multi-critères & 1~087~154 & 2~482,8 & 225,5 & 162,4 \\
Strict : bulletin ET cotisation & 616~327 & 1~889,7 & 192,0 & 125,0 \\
Élargi : bulletin OU cotisation & 902~080 & 2~285,3 & 215,6 & 149,9 \\
Diagnostic de non-réponse & 1~087~154 & 2~482,8 & 225,5 & 162,4 \\
\bottomrule
\end{tabular}
\begin{tablenotes}\footnotesize
\item \emph{Notes} : les quatre lignes utilisent la même rémunération
annualisée. Une non-réponse ne vaut jamais réponse positive. Le diagnostic
conserve le classement central et publie les valeurs manquantes séparément.
Les cotisations additionnent la CNPS salariale et la CMU simulée. Source :
calculs des auteurs à partir de l'EHCVM 2021.
\end{tablenotes}
\end{threeparttable}
\end{table}

Le décile supérieur acquitte 66,6~\% de l'IRPP et 51,5~\% des cotisations. Le
Kakwani de l'IRPP, $0,452$, est près de trois fois celui de la TVA centrale.
Conditionnellement aux salaires assujettis, l'IRPP représente 1,68~\% de la
masse salariale au premier décile, 5,43~\% au cinquième et 11,6~\% au dixième;
les cotisations correspondantes restent proches de 6,3--6,6~\%, en raison du
taux CNPS presque proportionnel et de la petite composante CMU. La figure
\ref{fig:direct} juxtapose cette incidence salariale à la charge rapportée au
revenu disponible de tous les ménages.
Cette comparaison ne doit pas être confondue avec l'effet agrégé. Dans le
scénario PDI central, le Gini passe de 0,3414 pour $Y_P$ à 0,3345 pour $Y_N$, soit une variation de
$-0,0069$. La cotisation retraite est alors de l'épargne et non un prélèvement
CEQ, même si sa charge statutaire reste décrite dans le tableau
\ref{tab:direct}. Dans le scénario PGT, le Gini passe de 0,3440 à 0,3336; 0,124~\%
de la population appartient toutefois à un ménage dont $Y_P$ reconstruit est
négatif. Le prélèvement est intensément progressif parmi les personnes exposées,
mais son champ ne couvre qu'une minorité des ménages.

La décomposition de la CMU éclaire aussi l'écart entre droit et observation. Le
tirage contraint par les réponses du module santé retient 50~641 participants
pauvres non contributifs et 885~822 non-pauvres, en dessous des cibles de
75~000 et 1~425~000 faute d'un vivier observable suffisant. Parmi les
non-pauvres, 537~231 adultes contribuent. La variante statutaire appliquée à
tous les adultes produirait 187,8 milliards de FCFA, contre 6,45 milliards dans
la simulation de participation observée; cet écart est une sensibilité de
couverture, non une recette attendue.

Le tableau~\ref{tab:direct-scenarios} dissipe l'inversion qui apparaissait dans
une version antérieure du calcul. Tous les scénarios utilisent désormais la
même assiette salariale. Le critère strict retient 616~327 personnes pondérées
et produit 192,0 milliards d'IRPP; le critère élargi en retient 902~080 et
produit 215,6 milliards, contre 1~087~154 personnes et 225,5 milliards dans le
scénario central. Les non-réponses ne sont jamais recodées comme un indice de
formalité. Leur diagnostic conserve donc la masse centrale et publie séparément
les effectifs dont les marqueurs sont manquants.

L'emploi secondaire ne représente que 1,60~\% de la masse salariale centrale.
Partager également les demi-parts d'enfants entre les assujettis, plutôt que les
attribuer au mieux rémunéré, augmente l'IRPP de 0,04~\% seulement.

\begin{figure}[t]
\centering
\includegraphics[width=0.88\textwidth]{fig16_direct_burden_decile.png}
\caption{Charge des impôts directs et cotisations par décile dans le scénario
central multi-critères. Le panneau de gauche rapporte les prélèvements au
revenu disponible de tous les ménages; celui de droite les rapporte à la masse
salariale des seuls assujettis. Source : calculs des auteurs à partir de l'EHCVM 2021.}
\label{fig:direct}
\end{figure}

\subsection{Transferts directs, filets sociaux et pensions}
\label{sec:res-transfers}

Les quatre transferts publics centraux totalisent 125,8 milliards de FCFA et
atteignent environ 1,00 million de ménages distincts pondérés. Le tableau
\ref{tab:transfers} sépare montants directement observés, imputations statutaires
et calages. Les 252,6 milliards de pensions sont présentés à part parce qu'ils
constituent du revenu différé en PDI et un transfert en PGT.

\begin{table}[t]
\centering
\begin{threeparttable}
\caption{Transferts directs et pensions observés ou imputés}
\label{tab:transfers}
\small
\begin{tabular}{lrrl}
\toprule
Composante & Ménages bénéficiaires & Masse & Traitement principal \\
           & (pondérés) & (mds FCFA) & \\
\midrule
PSSN                       & 191~767 &  27,108 & PMT et calage administratif \\
Bourses publiques          &  22~632 &   6,090 & montant EHCVM observé \\
Prestations familiales     & 614~165 &  84,324 & barème et éligibilité CNPS \\
Accidents du travail       & 751~661 &   8,276 & valeur attendue AT/MP \\
\textbf{Transferts directs} & \textbf{1~000~914} & \textbf{125,798} & union des bénéficiaires \\
Pensions contributives     & 215~283 & 252,603 & PDI central / PGT en sensibilité \\
\bottomrule
\end{tabular}
\begin{tablenotes}\footnotesize
\item \emph{Notes} : les nombres de bénéficiaires sont des ménages pondérés et
ne s'additionnent pas entre programmes. La ligne agrégée compte leur union. Les
totaux ne comptent ni les pensions alimentaires ni les sommes reçues d'autres
ménages, car ces flux circulent entre agents privés et ne sont pas versés par
une administration. AT/MP signifie
« accidents du travail et maladies professionnelles ». Source : calculs des auteurs à partir de l'EHCVM 2021.
\end{tablenotes}
\end{threeparttable}
\end{table}

Le calage PSSN sélectionne 191~767 ménages contre une cible de 192~000 et
reproduit exactement la masse de 27,108 milliards. Parmi les ménages imputés,
85,9~\% se situent sous le seuil de pauvreté, contre 37,7~\% des ménages qui
déclarent directement recevoir le programme monétaire. Le faible chevauchement
pondéré, 972 ménages, confirme que le PMT est une reconstruction de ciblage et
non le registre administratif.

La variante « déclarés d'abord » conserve les 24~989 déclarants, puis complète
la cible par le PMT. Elle sélectionne 192~544 ménages pondérés, maintient la
masse à 27,108 milliards et ramène la part de bénéficiaires pauvres à 79,5~\%.
Le profil reste favorable aux ménages modestes, mais moins que celui du PMT
central. Cet écart mesure directement la dépendance du diagnostic distributif à
l'absence d'appariement avec le registre social. Les autres programmes déclarés
restent documentés : aides alimentaires, travaux HIMO, soutien COVID et
prestations de santé ne sont pas transformés en espèces en l'absence de montant
identifiable.

L'indice de ciblage $G_{Y_D}-C_R$ vaut 0,092, avec un intervalle Rao--Wu
$[0,040;0,149]$. Positif, il indique que les transferts sont plus concentrés vers
le bas de la distribution que le revenu disponible, sans signifier que chaque
composante est pro-pauvre. La figure~\ref{fig:transfers} montre notamment que le
PSSN est concentré dans les premiers déciles, alors que prestations familiales
et accidents du travail suivent la couverture contributive.

\begin{figure}[t]
\centering
\includegraphics[width=0.88\textwidth]{fig18_transfers_decile.png}
\caption{Composition des transferts publics directs par décile. Le PSSN est
imputé et calé; les autres composantes suivent les règles décrites dans la
méthode. Source : calculs des auteurs à partir de l'EHCVM 2021.}
\label{fig:transfers}
\end{figure}

\subsection{Réductions de prix : un bénéfice concentré parmi les ménages raccordés}

Le scénario central attribue 25,15 milliards de FCFA de réductions de prix :
8,69 milliards pour l'électricité, calés exactement sur la référence ANARE-CI,
et 16,46 milliards pour la tranche sociale de l'eau. Aucun montant n'est
attribué aux carburants dans le scénario central. La borne haute atteint 95,1 milliards, mais
elle cumule volontairement trois hypothèses favorables --- 40~\% d'abonnés
électriques sociaux, une marge de 20~\% sur l'eau de revendeur et un écart de
80 FCFA par litre de carburant --- dont la dernière ne décrit pas directement
2021.

L'accès aux réseaux explique le profil distributif. Le premier décile reçoit
2,1~\% du montant central et le dixième 30,1~\%. Le coefficient de concentration
vaut 0,430 (IC à 95~\% : $[0,405;0,454]$), supérieur au Gini du revenu
disponible; l'indice de ciblage vers les ménages modestes est donc négatif,
$-0,097$ (IC : $[-0,121;-0,072]$). Une réduction réglementée du prix peut ainsi
être socialement motivée sans être progressive si les ménages modestes sont
moins souvent raccordés.

\begin{figure}[t]
\centering
\includegraphics[width=0.88\textwidth]{fig19_subsidies_decile.png}
\caption{Réductions de prix attribuées aux ménages par décile. Le scénario
central valorise l'électricité et l'eau du réseau, mais aucun soutien direct au
carburant. Source : calculs des auteurs à partir de l'EHCVM 2021.}
\label{fig:subsidies}
\end{figure}

\subsection{Éducation et santé : des services importants, mais non monétaires}

Les services publics d'éducation attribués aux élèves atteignent
1~354,87 milliards de FCFA avant paiements des familles et 1~301,62 milliards
après frais d'inscription et contributions scolaires. Le résultat net tombe à
1~179,5 milliards avec les seules dépenses courantes et à 1~161,6 milliards
lorsque toutes les dépenses scolaires déclarées sont retranchées. Dans le
scénario central, le coefficient de concentration vaut 0,170
(IC : $[0,092;0,262]$), nettement inférieur
au Gini du revenu disponible : le service est progressif relativement au revenu,
même si le dixième décile reçoit 16,1~\% de la dépense nette contre 7,0~\% pour
le premier.

\begin{figure}[t]
\centering
\includegraphics[width=0.88\textwidth]{fig20_education_decile.png}
\caption{Dépense publique d'éducation, paiements directs des familles et
bénéfice net par décile. Source : calculs des auteurs à partir de l'EHCVM 2021.}
\label{fig:education-inkind}
\end{figure}

Pour la santé, les 108,53 milliards du périmètre central sont entièrement
répartis; 81,85 milliards subsistent après les paiements directs aux structures
publiques. La borne budgétaire large donne 108,0 milliards nets. En revanche,
l'allocation fondée uniquement sur les consultations observées des trente
derniers jours ne donne que 13,9 milliards nets, car les frais déclarés sont
alors concentrés sur un petit nombre d'épisodes. Cet écart motive la moyenne de
recours de personnes comparables dans le scénario central. Le coefficient de
concentration est $-0,031$ (IC : $[-0,039;-0,023]$); le premier décile reçoit
11,5~\% du bénéfice net et le dixième 9,2~\%. Aucun de ces résultats n'utilise la
déclaration d'assurance maladie.

\begin{figure}[t]
\centering
\includegraphics[width=0.88\textwidth]{fig21_health_decile.png}
\caption{Dépense publique de santé, paiements directs et bénéfice net par
décile. Les soins ambulatoires centraux reposent sur le recours moyen par âge,
sexe et milieu; les hospitalisations sur l'usage observé. Source : calculs des auteurs à partir de l'EHCVM 2021.}
\label{fig:health-inkind}
\end{figure}

\subsection{Du revenu primaire au revenu final : baisse de l'inégalité, tension monétaire sur la pauvreté}

Le tableau~\ref{tab:income-complete} présente la séquence CEQ complète.
Le Gini passe de 0,3414 pour le revenu primaire PDI à 0,3336 pour le revenu
disponible, 0,3248 pour le revenu consommable et 0,3168 pour le revenu final.
Les impôts indirects réduisent donc l'inégalité relative tout en abaissant le
niveau monétaire : la pauvreté passe de 37,7~\% avant intervention à 42,0~\% au
revenu consommable. Après valorisation des services d'éducation et de santé,
l'indicateur comptable tombe à 33,9~\%. Cette dernière valeur ne signifie pas
que les ménages peuvent payer leur alimentation avec un service public; elle
sert à comparer la valeur distribuée par l'État au coût de production.

\begin{table}[t]
\centering
\begin{threeparttable}
\caption{Concepts de revenu CEQ complets, convention PDI}
\label{tab:income-complete}
\small
\begin{tabular}{lrrr}
\toprule
Concept & Moyenne annuelle & Gini & Pauvreté \\
        & (FCFA par personne) & & (\%) \\
\midrule
Revenu primaire      & 547~537 & 0,3414 & 37,70 \\
Revenu net de marché & 539~285 & 0,3345 & 37,93 \\
Revenu brut          & 551~856 & 0,3406 & 37,38 \\
Revenu disponible    & 543~604 & 0,3336 & 37,47 \\
Revenu consommable   & 504~447 & 0,3248 & 42,01 \\
Revenu final         & 551~305 & 0,3168 & 33,90 \\
\bottomrule
\end{tabular}
\begin{tablenotes}\footnotesize
\item \emph{Notes} : les revenus sont exprimés aux prix de l'enquête. Le revenu
final ajoute les services d'éducation et de santé à leur coût de production;
son taux de pauvreté est un indicateur comptable, non une mesure de liquidité.
Source : calculs des auteurs à partir de l'EHCVM 2021.
\end{tablenotes}
\end{threeparttable}
\end{table}

\begin{figure}[t]
\centering
\includegraphics[width=0.9\textwidth]{fig22_income_concepts.png}
\caption{Moyenne des concepts de revenu par décile. Les branches revenu net de
marché et revenu brut sont des présentations parallèles autour du revenu
primaire. Source : calculs des auteurs à partir de l'EHCVM 2021.}
\label{fig:income-complete}
\end{figure}

La décomposition de Shapley attribue la réduction totale de Gini de 0,0247 aux
impôts indirects pour 0,0084, aux prélèvements directs pour 0,0070, à
l'éducation pour 0,0068, à la santé pour 0,0018 et aux paiements publics directs
pour 0,0007. Les réductions de prix accroissent légèrement l'inégalité
($-0,00016$ de contribution à la réduction), en cohérence avec leur concentration
parmi les ménages raccordés. Les six contributions somment exactement à la
variation totale; l'incertitude de chacune est évaluée par rééchantillonnage
selon le plan de sondage.

\begin{figure}[t]
\centering
\includegraphics[width=0.84\textwidth]{fig23_shapley.png}
\caption{Contribution de Shapley des six groupes d'instruments à la réduction
du Gini. Une valeur positive réduit l'inégalité. Source : calculs des auteurs à partir de l'EHCVM 2021.}
\label{fig:shapley}
\end{figure}

La mesure non anonyme complète ce diagnostic relatif. Sous PDI, 37,47~\% de la
population subit une perte de ressources sous le seuil et 4,34~\% devient
nouvellement pauvre; les pertes sous le seuil totalisent 142,73 milliards de
FCFA, contre 9,30 milliards de gains. Cette incidence élevée reflète le fait
que beaucoup de ménages pauvres paient une taxe indirecte, même sans franchir
le seuil. Elle ne se confond donc pas avec l'augmentation de 4,31 points du taux
de pauvreté. Avec les seuls prélèvements et paiements directs, la population
appauvrie n'est que de 1,09~\%, tandis que les gains sous le seuil atteignent
31,03 milliards : l'essentiel de l'appauvrissement monétaire provient des impôts
sur la consommation.

\begin{figure}[t]
\centering
\includegraphics[width=0.84\textwidth]{fig24_fiscal_impoverishment.png}
\caption{Appauvrissement fiscal et gains sous le seuil pour des seuils compris
entre 0,5 et 2 fois le seuil national. Source : calculs des auteurs à partir de l'EHCVM 2021.}
\label{fig:fiscal-impoverishment}
\end{figure}

\subsection{Validation externe et comparaison avec l'ENV 2014}

Les comptes nationaux définitifs et l'Annuaire des statistiques économiques de
l'ANStat fournissent le cadre macroéconomique principal. En 2021, la
consommation finale des ménages atteint 26~754 milliards de FCFA, les
importations 9~211 milliards, les impôts et droits sur les produits nets de
subventions 3~074 milliards, la valeur ajoutée de l'enseignement 2~971 milliards
et celle de la santé et de l'action sociale 743 milliards
\citep{anstat2023,anstatcna2023}. Le modèle représente 16~035 milliards de
consommation des ménages, soit 59,9~\% de l'agrégat exhaustif. Ce taux de
couverture constitue le dénominateur naturel pour lire les autres écarts : une
enquête auprès des ménages ne doit pas retrouver les achats des entreprises, des
administrations et des non-résidents.

La TVA directe et non déductible supportée par les ménages atteint 907,3
milliards. Les recettes 2021 de TVA intérieure et de TVA en douane totalisent
1~116,6 milliards dans l'annuaire ANStat, soit un ratio de 81,3~\%. Cette
comparaison élargie complète le rapprochement entre la seule TVA directe S3
(692,9 milliards) et la TVA intérieure DGI (556,3 milliards). Les accises et
droits de douane contenus dans les achats des ménages atteignent 253,0
milliards; les recettes douanières hors TVA, beaucoup plus larges, atteignent
1~161,6 milliards. Après déduction des 25,15 milliards de réductions de prix,
les impôts indirects nets simulés représentent 1~135,2 milliards, contre
3~074 milliards pour l'ensemble des impôts sur les produits nets de subventions.

\begin{table}[t]
\centering
\begin{threeparttable}
\caption{Rapprochement des simulations avec les agrégats ANStat et budgétaires, 2021}
\label{tab:macro-validation}
\small
\begin{tabular}{lrrr}
\toprule
Agrégat & Simulation & Référence & Simulation / référence \\
        & \multicolumn{2}{c}{(milliards de FCFA)} & (\%) \\
\midrule
Consommation finale des ménages  & 16~034,6 & 26~754,0 & 59,9 \\
TVA directe et non déductible    &    907,3 &  1~116,6 & 81,3 \\
Accises, achats des ménages      &     67,9 &    292,9 & 23,2 \\
Droits de douane, ménages        &    185,1 &    555,1 & 33,3 \\
Impôts indirects nets des aides de prix & 1~135,2 & 3~074,0 & 36,9 \\
IRPP, scénario central           &    225,5 &    591,6 & 38,1 \\
PSSN                             &     27,1 &     27,1 & 100,0 \\
Aide au secteur électrique       &      8,69&     57,5 & 15,1 \\
Services d'éducation attribués   &  1~354,9 &  2~971,0 & 45,6 \\
Services de santé attribués      &    108,5 &    743,0 & 14,6 \\
Paiements publics directs        &    125,8 &    906,8 & 13,9 \\
Pensions contributives          &    252,6 &    180,9 & 139,6 \\
\bottomrule
\end{tabular}
\begin{tablenotes}\footnotesize
\item \emph{Notes} : les références ne sont égales au périmètre simulé que
pour le PSSN. La TVA additionne DGI et DGD; les références d'accises et de droits de
douane reprennent les postes comparables de la loi de règlement, hors redevance
statistique et hors TVA. Les références d'éducation et de santé sont les valeurs ajoutées
publiques et privées des branches, non les seuls budgets distribuables. Le poste
électrique du TOFE inclut notamment du gaz imputé; l'aide d'exploitation ANARE
de 8,69 milliards est, elle, retrouvée exactement. La référence de paiements
publics couvre toutes les subventions et autres transferts, et celle des pensions
la seule CGRAE. Sources : \citet{anstat2023,anstatcna2023,anare2021,civreglement2021,civperformance2021}.
\end{tablenotes}
\end{threeparttable}
\end{table}

Le PIB nominal définitif 2021 est de 39~821,3 milliards de FCFA selon l'ANStat.
La TVA totale supportée par les ménages représente ainsi 2,28~\% du PIB, l'IRPP
central 0,57~\% et les services d'éducation attribués 3,40~\%. Les chiffres
2022 et 2023 servent à vérifier la stabilité des
ordres de grandeur, mais ils ne remplacent aucun paramètre 2021. La consommation
finale des ménages passe de 26~754 milliards en 2021 à 29~510 en 2022 et
34~059 en 2023 dans les comptes définitifs. L'annuaire publié antérieurement
indiquait 32~721 milliards pour 2023 : cet écart de millésime statistique est
conservé dans les sources, mais la série définitive est retenue ici; les impôts
nets sur les produits passent de 3~074 à 3~276 puis
4~431 milliards. Les séries ultérieures servent donc de comparaison temporelle,
non de calage rétroactif.

Ces ratios ne sont pas des objectifs d'ajustement. Le revenu et les achats sont
sous-déclarés dans l'enquête; les comptes nationaux couvrent des agents absents
de l'EHCVM; les recettes sont comptabilisées au moment du recouvrement; et
certaines références réunissent plusieurs instruments. Pour chaque comparaison,
nous indiquons la définition, le périmètre et la source au lieu de choisir le
scénario qui reproduirait mécaniquement une recette. Le poste administratif
« impôts sur revenus et salaires » atteint ainsi 591,6 milliards, mais il est
plus large que l'IRPP des ménages simulé. De même, le Document de programmation
budgétaire et économique pluriannuelle projetait 268~947 cotisants publics à la
CGRAE en 2021 \citep{civdpbep2021}, alors que l'estimation fondée sur l'EHCVM réunit salariés
publics et privés.

La comparaison avec \citet{akim2020} apporte un second repère. Leur modèle
OpenFisca sur l'ENV 2014 repondère les revenus vers les comptes nationaux,
classe principalement les ménages par revenu et ne propage pas les taxes sur
consommations intermédiaires. L'ajout des accises, droits de douane et
transferts directs rapproche les périmètres d'instruments, mais notre
modèle utilise le revenu disponible de consommation par tête, ne recale pas
l'assiette salariale et ajoute la TVA enchâssée. Le tableau~\ref{tab:akim} doit
donc être lu comme une comparaison de diagnostics, non comme une évolution
temporelle à méthodologie constante.

\begin{table}[t]
\centering
\begin{threeparttable}
\caption{Repères ENV 2014 et EHCVM 2021}
\label{tab:akim}
\small
\begin{tabularx}{\textwidth}{lrrX}
\toprule
Indicateur & ENV 2014 & EHCVM 2021 & Différence de champ \\
\midrule
IRPP simulé (mds FCFA) & 326,9 & 225,5 &
ENV recalée; EHCVM non recalée, scénario central \\
Part IRPP du D10 (\%) & 86,0 & 66,6 &
Concentration de l'impôt simulé \\
Part TVA du D1 (\%) & 5,0 & 1,95 &
TVA finale directe contre TVA directe + enchâssée S3 \\
Part TVA du D10 (\%) & 28,0 & 33,2 &
Classement par revenu contre $Y_D$ par tête \\
\bottomrule
\end{tabularx}
\begin{tablenotes}\footnotesize
\item \emph{Notes} : les valeurs ENV sont reprises de \citet{akim2020}; les
valeurs EHCVM sont calculées dans la présente étude. Les parts TVA sont des
parts de la masse, non des taux effectifs.
\end{tablenotes}
\end{threeparttable}
\end{table}

%=============================================================================
\section{Simulations de réformes}
\label{sec:reformes}

\subsection{Taxer les intrants avicoles : une charge concentrée en haut de la distribution}

La première simulation examine la suppression de l'exonération de TVA sur les
intrants de la filière avicole (provende, poussins d'un jour), avec
transmission partielle aux prix du poulet de chair. La clé du résultat est la
courbe d'Engel du poulet : sa part budgétaire croît de 0,06~\% de la
consommation au premier décile à 0,93~\% au dixième. La charge additionnelle
épouse ce profil --- le décile supérieur en supporterait 45~\% --- et l'indice
de Kakwani de la mesure ressort à $+0,24$ : la charge est progressive. L'effet
sur la pauvreté est faible dans ce modèle statique ($+0,01$ point de $P_0$,
moins de 500 ménages basculant sous le seuil dans le scénario de transmission
maximale). Dans les hypothèses retenues, une taxation des intrants avicoles
fait donc
partie des élargissements d'assiette dont le coût distributif simulé est le plus
faible.

\subsection{Taxer la vente finale à 9~\% à TVA incorporée constante}

La seconde simulation applique une taxe de 9~\% à la vente finale de produits
initialement à taux nul. Elle maintient le vecteur $c$ de TVA incorporée
constant et ne simule donc pas une reclassification juridique complète, qui
ouvrirait des droits à déduction sur les intrants. Nous distinguons les
produits agricoles et l'ensemble large hors services publics structurellement
exonérés. La catégorie du commerce ne contient aucun produit éligible.

Sous S3, la variante agricole collecte 54,1 milliards de FCFA et la variante
large 97,2 milliards. Sans usage des recettes, elles augmentent respectivement
$P_0$ de 0,26 et 0,47 point. Les intervalles Rao--Wu excluent zéro. La variante
large fait basculer 141~356 personnes sous le seuil; la variante agricole,
75~938. Ces résultats mesurent le choc de prix brut, pas le bien-être associé à
une politique complète. Ils correspondent respectivement à 1~455 et 1~404
nouveaux pauvres par milliard de FCFA de recette brute dans ces deux
contrefactuels S3; ce ratio disparaît comme critère de bien-être une fois les
recettes recyclées.

\begin{table}[t]
\centering
\begin{threeparttable}
\caption{Taxe de 9~\% sur la vente finale et recyclage universel, scénario S3}
\label{tab:reforme}
\begin{tabular}{lrrrr}
\toprule
Ensemble & Recette & Transfert & $\Delta P_0$ brut & $\Delta P_0$ recyclé \\
         & (mds FCFA) & (FCFA/personne) & \multicolumn{2}{c}{(point)} \\
\midrule
Agricole & 54,1 & 1~815 & 0,26 [0,14 ; 0,38] & $-0,15$ [$-0,24$ ; $-0,06$] \\
Large    & 97,2 & 3~260 & 0,47 [0,31 ; 0,64] & $-0,23$ [$-0,35$ ; $-0,12$] \\
\bottomrule
\end{tabular}
\begin{tablenotes}\footnotesize
\item \emph{Notes} : S3 fait varier la taxation effective par COICOP et
décile. Le transfert est uniforme par personne et son coût agrégé égale
exactement la recette nominale simulée. La recette et le transfert sont
réestimés dans chacune des 500 réplications Rao--Wu. Source : calculs des auteurs à partir de l'EHCVM 2021.
\end{tablenotes}
\end{threeparttable}
\end{table}

Le recyclage inverse le diagnostic. Le transfert agricole réduit $P_0$ de 0,15
point; le transfert associé à la réforme large le réduit de 0,23 point. Dans
cette dernière variante, aucun ménage initialement non pauvre ne passe sous le
seuil après impôt et transfert, tandis qu'environ 66~980 personnes initialement
pauvres le franchissent vers le haut. Pour l'agriculture, 1~982 personnes
deviennent pauvres mais 45~115 en sortent. La redistribution universelle est favorable aux ménages modestes
parce que chaque personne reçoit le même montant, alors que la taxe augmente
avec la dépense formelle.

Ce résultat est cohérent avec l'analyse de \citet{warwick2022} : juger une exonération
à partir de son seul effet sur les prix ignore l'instrument qui pourrait être
financé par sa suppression. Il ne suffit cependant pas à recommander la réforme.
Le transfert universel suppose une capacité de paiement, ne comporte aucun coût
administratif et ne modélise ni ajustement des quantités ni effet sur la
production. Il constitue une référence transparente contre laquelle un ciblage
réaliste pourra être comparé.

\begin{figure}[t]
\centering
\includegraphics[width=0.86\textwidth]{fig_reform_fgt_impact.png}
\caption{Variation de la pauvreté après taxation à 9~\% de la vente finale,
avant et après recyclage universel. Le vecteur de TVA incorporée reste constant;
barres d'erreur Rao--Wu à 95~\%.
Source : calculs des auteurs à partir de l'EHCVM 2021.}
\label{fig:reform-recycling}
\end{figure}

%=============================================================================
\clearpage
\section{Portée et limites}
\label{sec:limites}

Les résultats mesurent une incidence comptable sur les ménages observés en
2021. Cinq limites en définissent la portée. Elles n'annulent pas les écarts
mis en évidence, mais déterminent les interprétations qui peuvent leur être
données. L'annexe~\ref{annexe:limites} détaille, pour chaque point, le biais
possible, les contrôles disponibles et les données qui permettraient de le
réduire.

\paragraph{Mesure des ressources et couverture administrative.}
La consommation de l'EHCVM sert d'approximation du revenu disponible parce
qu'elle est mieux mesurée que le revenu courant, mais elle n'est pas un revenu
de marché. Le revenu primaire est donc reconstruit plutôt qu'observé. Les
salaires sont sous-déclarés, les revenus indépendants, fonciers et mobiliers
restent incomplets, et l'enquête n'est appariée ni aux fichiers fiscaux ni aux
registres sociaux. La déclaration d'assurance maladie est particulièrement
fragile : elle ne retrouve qu'une partie des effectifs administratifs de la
CMU. Les conventions PDI et PGT, les variantes d'assujettissement et la
publication des écarts de couverture rendent ces incertitudes visibles. Elles
ne remplacent pas un appariement individuel, qui permettrait surtout de mieux
mesurer les prélèvements directs et l'accès effectif aux programmes.

\paragraph{Incidence statique et taxation effective.}
Le modèle maintient fixes les quantités achetées, les coefficients techniques
et l'offre. Il décrit ainsi la transmission comptable de court terme, sans
substitution entre produits ou intrants, réaction de production, formalisation
ni effet d'équilibre général. De plus, l'EHCVM ne renseigne pas
systématiquement le statut fiscal du vendeur. Le scénario S3 impute donc un
coefficient de taxation à partir du produit et du niveau de vie. Sa matrice est
publiée intégralement, mais elle reste exogène et n'est calée sur aucune cible
ivoirienne. La variante S2 et la transmission intégrale encadrent cette
hypothèse, mais ne constituent pas des réponses comportementales. Des données de caisse, de
facturation électronique et de contrôle fiscal permettraient d'estimer la
collecte effective; un modèle de demande et de production permettrait ensuite
d'évaluer l'ajustement des ménages et des entreprises.

\paragraph{Matrice de production et fiscalité indirecte.}
Le TRE utilisé date de 2023, soit deux ans après l'enquête. La concordance vers
ses 48 produits couvre 92,5~\% de la consommation pondérée, contre 95,2~\%
pour ICIO, et applique une même part domestique aux différents usages d'un
produit. Les accises et les droits de douane sont en outre associés à des
postes de consommation plus agrégés que les lignes du Système harmonisé. Leur
transmission directe est simulée, mais pas encore leur propagation par les
intrants; leur charge totale constitue donc vraisemblablement une borne basse.
Les variantes à prix courants, à prix constants et avec ICIO confirment la
stabilité au choix de matrice. En revanche, la sensibilité juridique du droit à
déduction réduit fortement la TVA incorporée : ce paramètre reste donc une
source majeure d'incertitude. Un TRE 2021 distinguant les importations selon
leur usage, une nomenclature administrative des droits à déduction et une
concordance douanière plus fine réduiraient cette incertitude.

\paragraph{Allocation des paiements et des services publics.}
Le PMT reproduit une règle de ciblage du PSSN sans identifier les bénéficiaires
du registre. Les prestations familiales et les accidents du travail reposent
sur des règles statutaires et des indicateurs de couverture. Les réductions de
prix de l'électricité et de l'eau sont déduites des factures et des barèmes;
faute de prix de parité 2021 comparable, aucun soutien au carburant n'entre dans
le scénario central. Enfin, l'éducation et la santé sont valorisées à leur coût
moyen, sans mesurer complètement la qualité, l'assiduité ou les différences de
coût entre établissements. Ces conventions peuvent modifier la répartition de
chaque programme, mais pas la distinction essentielle entre un paiement
monétaire et un service en nature. Les registres PSSN, les catégories
tarifaires inscrites sur les factures et des comptes sectoriels plus détaillés
permettraient de tester ces allocations.

\paragraph{Comparabilité et portée temporelle.}
La consommation mesurée dans l'EHCVM représente 59,9~\% de la consommation
finale des ménages des comptes nationaux. Les comptes et les recettes fiscales
couvrent en outre des entreprises et des administrations absentes de l'enquête.
Forcer les masses simulées à reproduire ces agrégats imposerait des hypothèses
fortes sur leur répartition entre ménages. L'étude publie donc les écarts de
périmètre et utilise les valeurs 2022--2023 comme contrôles, non comme cibles de
calage. Les résultats décrivent l'incidence sur les ménages ivoiriens en 2021;
ils ne constituent ni une réconciliation exhaustive des recettes publiques ni
une prévision de l'effet d'une réforme à une autre date.
\clearpage
\section{Conclusion}
\label{sec:conclusion}

En 2021, l'ensemble des impôts et des dépenses publiques réduit l'inégalité
mesurée en Côte d'Ivoire, mais ne protège pas tous les ménages contre la perte
de pouvoir d'achat. Dans le scénario central, le Gini passe de 0,3414 pour le
revenu primaire à 0,3168 pour le revenu final. Au revenu consommable, avant la
valorisation de l'éducation et de la santé, le taux de pauvreté atteint pourtant
42,0~\%, contre 37,7~\% au revenu primaire, et 4,34~\% de la population passe
sous le seuil national. Les pertes monétaires sous ce seuil atteignent
142,7 milliards de FCFA. La baisse de l'inégalité et l'appauvrissement fiscal
mesurent donc deux dimensions distinctes du même système.

La composition des interventions explique ce contraste. L'impôt sur le revenu
est très progressif, mais son assiette ne touche qu'une minorité de ménages. Les
paiements publics directs restent trop faibles pour compenser la charge des
taxes sur la consommation : les quatre programmes centraux totalisent
125,8 milliards de FCFA, dont 27,1 milliards pour le PSSN. L'éducation et la
santé fournissent l'essentiel du bénéfice redistributif, avec respectivement
1~301,6 et 81,85 milliards nets des paiements des familles. Ces services
améliorent le revenu final au sens comptable, mais ne financent pas les achats
courants. Le taux de pauvreté de 33,9~\% calculé après leur valorisation ne doit
donc pas être lu comme une mesure de liquidité.

La simulation livre une implication directe pour la politique fiscale. Une
taxe de 9~\% sur la vente finale de certains produits accroît la pauvreté de
0,47 point lorsque la recette n'est pas restituée. Ce calcul maintient la TVA
incorporée constante et ne représente donc pas une reclassification juridique
complète. Un transfert uniforme financé par les 97,2 milliards de FCFA
collectés réduit au contraire la pauvreté de 0,23 point.
L'élargissement de l'assiette et la protection des ménages modestes peuvent donc
être compatibles, mais seulement si le mécanisme de restitution est explicite,
financé et effectivement accessible. Le transfert universel simulé sert ici de
référence budgétairement neutre; il ne tranche pas entre universalité et ciblage.

Les prochaines améliorations doivent réduire les écarts entre simulation et
administration plutôt qu'ajouter des conventions invisibles. L'appariement de
l'EHCVM avec les registres fiscaux et sociaux permettrait d'observer la taxation
effective, l'assujettissement salarial et les bénéficiaires du PSSN. Un TRE de
2021 distinguant les importations selon leur usage améliorerait la mesure des
taxes incorporées aux prix. Enfin, un modèle de demande et de production
permettrait d'étudier les changements de consommation, de formalisation et de
coûts administratifs. Ces extensions rendraient possible la comparaison de
réformes à la fois redistributives, finançables et applicables.
\clearpage
\bibliographystyle{apalike}
\bibliography{references}

%=============================================================================
\clearpage
\appendix

\section{Barèmes des impôts directs (CGI 2021)}
\label{annexe:baremes}

L'impôt sur les salaires (IS) vaut 1,5~\% de la base imposable, définie comme
80~\% du salaire brut imposable (SBI), soit 1,2~\% du SBI. La contribution
nationale (CN) s'applique à la même base selon les tranches mensuelles
suivantes : 0~\% jusqu'à 50~000 FCFA ; 1,5~\% de 50~001 à 130~000 ; 5~\% de
130~001 à 200~000 ; 10~\% au-delà. La base de l'IGR vaut
$B_{\mathrm{IGR}} = 0{,}8\,(SBI - IS - CN)$ ; le quotient
$Q = B_{\mathrm{IGR}}/N$ est soumis au barème
mensuel du tableau~\ref{tab:bareme-igr}, l'impôt par part étant ensuite
multiplié par $N$ puis réduit de 15~\%. Le nombre de parts $N$ (plafonné à 5)
vaut 1 pour un célibataire, divorcé ou séparé, 1,5 pour un veuf, 2 pour un
marié, plus 0,5 par enfant à charge. La cotisation retraite CNPS salariale vaut
6,3~\% du SBI plafonné à 1~647~315 FCFA par mois; la contribution CMU simulée
vaut 1~000 FCFA par mois et par adulte contributif participant.

\begin{table}[h]
\centering
\begin{threeparttable}
\caption{Barème mensuel de l'IGR par part ($Q = B_{\mathrm{IGR}}/N$, en FCFA)}
\label{tab:bareme-igr}
\begin{tabular}{lc}
\toprule
Tranche de $Q$ & Impôt par part \\
\midrule
$Q \le 25~000$            & $0$ \\
$25~000 < Q \le 45~000$   & $Q \times 10/110 - 2~273$ \\
$45~000 < Q \le 60~000$   & $Q \times 15/115 - 4~076$ \\
$60~000 < Q \le 80~000$   & $Q \times 20/120 - 7~083$ \\
$80~000 < Q \le 100~000$  & $Q \times 25/125 - 11~000$ \\
$100~000 < Q \le 120~000$ & $Q \times 30/130 - 15~769$ \\
$120~000 < Q \le 140~000$ & $Q \times 35/135 - 21~296$ \\
$140~000 < Q \le 160~000$ & $Q \times 40/140 - 27~500$ \\
$160~000 < Q \le 180~000$ & $Q \times 45/145 - 34~310$ \\
$180~000 < Q \le 200~000$ & $Q \times 50/150 - 41~667$ \\
$Q > 200~000$             & $Q \times 60/160 - 57~813$ \\
\bottomrule
\end{tabular}
\begin{tablenotes}\footnotesize
\item \emph{Notes} : IGR final $= (\text{impôt par part} \times N) \times 0{,}85$.
Source : Code général des impôts de Côte d'Ivoire, édition 2021.
\end{tablenotes}
\end{threeparttable}
\end{table}

\section{Concordance de la nomenclature intermédiaire vers le TRE}
\label{annexe:mapping-io}

Le matériel de réplication du projet contient une concordance détaillée,
initialement fournie dans le programme Stata de référence, qui associe chaque
poste de consommation à l'un des 52 secteurs intermédiaires. Le pipeline R
extrait ce tableau, vérifie l'unicité de chaque affectation, puis applique la
passerelle ci-dessous. Les doublons apparents correspondent à des catégories
intermédiaires plus détaillées que le TRE : par exemple, les secteurs 1 et 2
sont tous deux regroupés dans \texttt{A01}. Aucun item n'est réparti entre
plusieurs branches. La liste exhaustive au niveau item est exportée dans
\texttt{13\_04\_itemid\_to\_tre\_mapping\_current.xlsx}; la passerelle est
exportée dans \texttt{13\_05\_wb52\_to\_tre48\_bridge\_current.xlsx}.

\begin{longtable}{rr@{\hspace{1.6cm}}rr}
\caption{Passerelle des 52 secteurs intermédiaires vers les 48 produits du TRE}
\label{tab:mapping-52-48}\\
\toprule
Secteur interm. & Code TRE & Secteur interm. & Code TRE \\
\midrule
\endfirsthead
\toprule
Secteur interm. & Code TRE & Secteur interm. & Code TRE \\
\midrule
\endhead
1  & \texttt{A01} & 27 & \texttt{C23} \\
2  & \texttt{A01} & 28 & \texttt{C24} \\
3  & \texttt{A02} & 29 & \texttt{C25} \\
4  & \texttt{A03} & 30 & \texttt{C26} \\
5  & \texttt{A04} & 31 & \texttt{C27} \\
6  & \texttt{A05} & 32 & \texttt{C28} \\
7  & \texttt{A06} & 33 & \texttt{C29} \\
8  & \texttt{B07} & 34 & \texttt{C30} \\
9  & \texttt{C08} & 35 & \texttt{D31} \\
10 & \texttt{C09} & 36 & \texttt{E32} \\
11 & \texttt{C10} & 37 & \texttt{F33} \\
12 & \texttt{C11} & 38 & \texttt{G34} \\
13 & \texttt{C12} & 39 & \texttt{H35} \\
14 & \texttt{C13} & 40 & \texttt{I36} \\
15 & \texttt{C13} & 41 & \texttt{J37} \\
16 & \texttt{C14} & 42 & \texttt{K38} \\
17 & \texttt{C15} & 43 & \texttt{L39} \\
18 & \texttt{C16} & 44 & \texttt{M40} \\
19 & \texttt{C17} & 45 & \texttt{N41} \\
20 & \texttt{C18} & 46 & \texttt{O42} \\
21 & \texttt{C19} & 47 & \texttt{P43} \\
22 & \texttt{C19} & 48 & \texttt{Q44} \\
23 & \texttt{C20} & 49 & \texttt{R45} \\
24 & \texttt{C21} & 50 & \texttt{S46} \\
25 & \texttt{C21} & 51 & \texttt{T47} \\
26 & \texttt{C22} & 52 & \texttt{U48} \\
\bottomrule
\end{longtable}

\clearpage
\section{Paramètres publiés du scénario S3}
\label{annexe:s3-parameters}

Le scénario S3 applique la règle de l'équation~\eqref{eq:alpha-s3}. Le tableau
\ref{tab:s3-matrix} publie les paramètres et chaque coefficient utilisé par le
code. Les divisions 98 et 99 désignent respectivement la valeur d'usage des
biens durables et les postes hors champ; elles ne supportent aucune TVA. Ces
valeurs sont exogènes. Aucune cible de recette, de pauvreté ou de formalité
ivoirienne n'a servi à les calibrer.

\begin{table}[h]
\centering
\caption{Matrice exacte des coefficients de taxation effective du scénario S3}
\label{tab:s3-matrix}
\scriptsize
\setlength{\tabcolsep}{2.6pt}
\resizebox{\textwidth}{!}{%
\begin{tabular}{lrrrrrrrrrrrr}
\toprule
Fonction COICOP & $a_k$ & $s_k$ & D1 & D2 & D3 & D4 & D5 & D6 & D7 & D8 & D9 & D10 \\
\midrule
1 Alimentation & .120 & .034 & .120 & .154 & .188 & .222 & .256 & .290 & .324 & .358 & .392 & .426 \\
2 Alcool et tabac & .480 & .024 & .480 & .504 & .528 & .552 & .576 & .600 & .624 & .648 & .672 & .696 \\
3 Vêtements et chaussures & .220 & .030 & .220 & .250 & .280 & .310 & .340 & .370 & .400 & .430 & .460 & .490 \\
4 Logement et énergie & .620 & .022 & .620 & .642 & .664 & .686 & .708 & .730 & .752 & .774 & .796 & .818 \\
5 Équipement du logement & .220 & .028 & .220 & .248 & .276 & .304 & .332 & .360 & .388 & .416 & .444 & .472 \\
6 Santé & .300 & .040 & .300 & .340 & .380 & .420 & .460 & .500 & .540 & .580 & .620 & .660 \\
7 Transport & .250 & .038 & .250 & .288 & .326 & .364 & .402 & .440 & .478 & .516 & .554 & .592 \\
8 Information et communication & .780 & .015 & .780 & .795 & .810 & .825 & .840 & .855 & .870 & .885 & .900 & .915 \\
9 Loisirs et culture & .280 & .032 & .280 & .312 & .344 & .376 & .408 & .440 & .472 & .504 & .536 & .568 \\
10 Éducation & .550 & .025 & .550 & .575 & .600 & .625 & .650 & .675 & .700 & .725 & .750 & .775 \\
11 Restaurants et hôtels & .120 & .038 & .120 & .158 & .196 & .234 & .272 & .310 & .348 & .386 & .424 & .462 \\
12 Assurance et finance & .850 & .010 & .850 & .860 & .870 & .880 & .890 & .900 & .910 & .920 & .930 & .940 \\
13 Soins personnels et divers & .150 & .034 & .150 & .184 & .218 & .252 & .286 & .320 & .354 & .388 & .422 & .456 \\
98 Valeur d'usage des durables & .000 & .000 & .000 & .000 & .000 & .000 & .000 & .000 & .000 & .000 & .000 & .000 \\
99 Hors champ & .000 & .000 & .000 & .000 & .000 & .000 & .000 & .000 & .000 & .000 & .000 & .000 \\
\bottomrule
\end{tabular}}
\begin{minipage}{0.98\textwidth}
\footnotesize\emph{Notes} : $a_k$ est le coefficient au premier décile et
$s_k$ sa hausse par décile. D1--D10 donnent directement $\alpha_{kd}$. Le
classeur de réplication \texttt{06\_03\_vat\_informality\_parameters.xlsx}
publie aussi les paramètres S2 et la matrice sous forme lisible par machine.
\end{minipage}
\end{table}

Après pondération par les dépenses observées de chaque décile, le profil
implicite de taxation effective vaut 0,266 en D1, 0,294 en D2, 0,325 en D3,
0,354 en D4, 0,386 en D5, 0,421 en D6, 0,450 en D7, 0,486 en D8, 0,522 en D9
et 0,578 en D10. Ce profil est une conséquence de la matrice et des paniers de
consommation; il n'est pas une cible de calage. S3 doit donc être lu comme une
hypothèse structurante publiée, encadrée par S2 et par la transmission intégrale.
\clearpage
\section{Détail des limites et données nécessaires}
\label{annexe:limites}

Cette annexe conserve les réserves techniques dans un format opérationnel.
Chaque point distingue la source d'incertitude, son effet possible, le contrôle
mis en œuvre et l'information qui permettrait de le réduire.

\begin{enumerate}
\item \textbf{Point d'entrée et pensions.} La consommation observée est
l'ancre du revenu disponible; le revenu primaire est reconstruit. Cette
convention peut modifier le niveau des revenus antérieurs sans changer le
revenu disponible observé. Les conventions PDI et PGT encadrent le traitement
des pensions. Des données conjointes de revenu, consommation et pensions
administratives permettraient une reconstruction directe.

\item \textbf{Réactions comportementales.} Les quantités, l'offre et les
coefficients techniques sont fixes. Les simulations mesurent donc un effet
comptable de court terme, qui peut différer de l'effet après substitution ou
formalisation. Les scénarios de transmission testent l'incidence immédiate; un
système de demande et un modèle de production seraient nécessaires pour mesurer
les ajustements.

\item \textbf{Informalité des achats.} Le statut fiscal du vendeur n'est pas
observé pour chaque transaction. Les scénarios S2, S3 et de transmission
intégrale encadrent la charge, sans identifier la probabilité réelle de
collecte. Des tickets de caisse, des factures électroniques ou un appariement
avec les déclarations de TVA permettraient de l'estimer.

\item \textbf{TRE et concordance sectorielle.} Le TRE 2023 est postérieur à
l'EHCVM 2021, sa concordance couvre 92,5~\% de la consommation pondérée et la
part importée est commune aux différents usages d'un produit. Les matrices à
prix constants et ICIO servent de contrôles. Un TRE 2021 séparant les usages
domestiques et importés serait préférable.

\item \textbf{Impôts directs et CMU.} L'assujettissement salarial est imputé et
les revenus non salariaux sont incomplets. Sur une assiette salariale commune,
l'IRPP vaut 192,0 milliards dans le scénario strict, 215,6 milliards dans le
scénario élargi et 225,5 milliards dans le scénario central. Une non-réponse
n'est jamais recodée comme emploi formel. La déclaration d'assurance ne retrouve
que 50~641 participants pauvres et 885~822 non-pauvres, contre des cibles de
75~000 et 1~425~000. Les fichiers fiscaux, sociaux et de la CMU permettraient de
remplacer ces bornes par une observation administrative.

\item \textbf{Accises et droits de douane.} Les taux sont affectés à des postes
plus agrégés que les lignes douanières et les parts importées proviennent du
TRE. La transmission directe est simulée, mais pas la propagation des taxes
payées sur les intrants. Les estimations sont donc susceptibles de sous-estimer
la charge totale. Une concordance entre dépenses, lignes du Système harmonisé
et usages du TRE permettrait de compléter la propagation.

\item \textbf{Paiements publics et PSSN.} Le PMT reproduit un ciblage mais ne
reconnaît pas les bénéficiaires réels. Le chevauchement avec les déclarations de
l'enquête et la variante qui retient d'abord tous les déclarants sont publiés.
Les prestations familiales et les
accidents du travail suivent des règles statutaires; les aides en nature ne sont
pas monétisées sans montant fiable. L'appariement au registre social et aux
fichiers des organismes payeurs est le contrôle décisif.

\item \textbf{Réductions publiques de prix.} L'électricité est calée sur
8,69 milliards de FCFA d'aide d'exploitation, mais les abonnés sociaux sont
identifiés par leur rang de facture. Pour l'eau, le volume est déduit du barème
et le prix des revendeurs n'est pas observé. Aucun soutien au carburant n'est
retenu au centre faute de prix de parité 2021 auditable. Les catégories
contractuelles, volumes facturés et prix de référence officiels permettraient
de tester ces imputations.

\item \textbf{Éducation et santé.} Les services sont valorisés à un coût moyen
sans mesurer intégralement leur qualité, leur assiduité ou les écarts entre
établissements. La santé lisse les consultations ambulatoires par âge, sexe et
milieu; une variante fondée sur le seul usage observé donne une borne plus
basse. Le coût attribué n'est ni un revenu liquide ni un consentement à payer.
Des comptes sectoriels par niveau de service et des données administratives
d'utilisation affineraient l'incidence.

\item \textbf{Comparaison avec les comptes nationaux.} La consommation de
l'EHCVM ne représente que 59,9~\% de l'agrégat des comptes nationaux, tandis que
les recettes couvrent aussi les entreprises et les administrations. L'étude
publie ces écarts et les séries 2022--2023 plutôt que de caler globalement les
ménages. Une réconciliation complète exigerait des comptes de ménages par
niveau de vie et des recettes ventilées par secteur institutionnel.
\end{enumerate}
\section{Accès au code et protocole de réplication}
\label{annexe:pipeline}

\subsection{Dépôt public et version de référence}

Le code et la documentation méthodologique sont maintenus dans le dépôt GitHub
public \url{https://github.com/cae-ins/anstat-microsim}. Le dépôt peut être
consulté et cloné sans compte GitHub. La branche par défaut est
\texttt{main}; la branche \texttt{rewrite-r} porte l'implémentation R utilisée
dans cette étude. Le clonage de cette branche s'effectue depuis un terminal par
les commandes suivantes :

\begin{verbatim}
git clone --branch rewrite-r --single-branch \
  https://github.com/cae-ins/anstat-microsim.git
cd anstat-microsim
\end{verbatim}

Une branche est un objet évolutif. Pour qu'une réplication puisse identifier
sans ambiguïté la version du modèle, le numéro de révision obtenu par
\texttt{git rev-parse HEAD} doit donc être conservé avec les résultats. Lors de
la diffusion définitive du papier, cette révision sera gelée dans une
\emph{release} GitHub portant le même numéro de version que le manuscrit. Cette
précaution distingue l'adresse durable du projet de l'état exact du code ayant
produit un tableau ou une figure.

\subsection{Accès aux microdonnées et périmètre du dépôt}

Le dépôt public fournit le code et la documentation, mais ne duplique pas les
microdonnées. L'EHCVM 2021--2022 peut être obtenue par deux canaux officiels :
le centre de calcul de l'ANStat
(\url{https://centredecalcul.anstat.ci/index.php/welcome}) et le catalogue de
microdonnées de la Banque mondiale
(\url{https://microdata.worldbank.org/index.php/catalog/6273}), où l'enquête
porte la référence \texttt{CIV\_2021\_EHCVM-2\_v01\_M}. Le catalogue fournit
les données anonymisées de diffusion publique, leur dictionnaire et les
questionnaires. Dans les deux cas, l'utilisateur doit respecter les conditions
d'accès, de confidentialité et de citation affichées au moment du
téléchargement, et consigner la source ainsi que la date de téléchargement.

Les fichiers individuels de l'EHCVM, les données intermédiaires et les sorties
générées sont exclus du versionnement Git. Cette séparation évite de
redistribuer des microdonnées en dehors des conditions acceptées par
l'utilisateur et oblige le réplicateur à acquérir les fichiers auprès d'un des
catalogues officiels. Le tableau~\ref{tab:acces-replication} précise cette
frontière.

\begin{longtable}{>{\RaggedRight\arraybackslash}p{3.2cm}
  >{\RaggedRight\arraybackslash}p{5.0cm}
  >{\RaggedRight\arraybackslash}p{6.1cm}}
\caption{Accessibilité des composantes de la réplication}
\label{tab:acces-replication}\\
\toprule
Composante & Emplacement attendu & Accessibilité et rôle \\
\midrule
\endfirsthead
\toprule
Composante & Emplacement attendu & Accessibilité et rôle \\
\midrule
\endhead
Code R et documentation &
\path{05_scripts_R/} et \path{00_documentation/} &
Publics sur GitHub; ils définissent les transformations, les hypothèses,
les contrôles et la production du papier. \\
Microdonnées EHCVM harmonisées &
\path{01_data_sources/Dataout/} &
Accessibles par les catalogues de l'ANStat et de la Banque mondiale, puis à
placer localement dans ce répertoire; elles ne sont pas recopiées dans Git. \\
Modules ménages détaillés &
\path{01_data_sources/Datain/Menage/} &
Accessibles dans la livraison EHCVM. Les modules d'emploi et de couverture
maladie sont notamment requis pour les prélèvements directs et la CMU. \\
Paramètres fiscaux et concordances &
\path{01_data_sources/} &
Le classeur de taux par poste de consommation et les concordances sectorielles
doivent être placés aux chemins lus par les scripts. Leur version doit être
identique à celle associée à la release de réplication. \\
Matrices entrées--sorties &
\path{01_data_sources/IO/} &
Les fichiers ICIO et les TRE ivoiriens courant et constant sont à obtenir
auprès de leurs producteurs puis à déposer dans ce répertoire. \\
Intermédiaires et résultats &
\path{02_data_intermediate/} et \path{07_reports/} &
Non versionnés car entièrement reconstruits. Les fichiers parquet assurent les
passages entre étapes; les classeurs et figures constituent les sorties
publiables. \\
\bottomrule
\end{longtable}

Les noms techniques sont utiles ici parce qu'ils permettent au réplicateur de
placer les entrées sans ambiguïté. Les trois fichiers EHCVM harmonisés attendus
par le cœur du pipeline sont
\path{ehcvm_conso_CIV2021.dta},
\path{ehcvm_welfare_2b_CIV2021.dta} et
\path{ehcvm_individu_CIV2021.dta}. Les prélèvements directs mobilisent les
modules d'emploi et de santé; les transferts ajoutent les modules d'éducation,
d'emploi, de pensions, de transferts privés et de programmes sociaux. L'incidence
indirecte requiert le classeur de TVA, les paramètres 2021 d'accises et du TEC,
la concordance sectorielle, ICIO 2020 et les deux TRE 2023. La présence d'un nom
de fichier dans le code ne vaut pas autorisation de redistribution de son
contenu.

Ces trois fichiers harmonisés ne résultent pas d'un simple changement de nom
des fichiers téléchargés. Ils constituent la couche préparée officielle :
les modules individuels et ménages sont d'abord assemblés, les dépenses par
poste sont annualisées et agrégées, puis l'agrégat de bien-être et son
déflateur spatial sont construits selon la variante nationale 2b. Les
programmes de préparation correspondants sont conservés dans
\path{01_data_sources/Programs/}; ils produisent successivement les fichiers
d'individus, de consommation et de bien-être attendus dans
\path{01_data_sources/Dataout/}. Un utilisateur partant des modules bruts doit
donc exécuter cette préparation, après avoir adapté uniquement les chemins
d'entrée et de sortie à son installation. Cette étape préalable requiert Stata;
elle n'est pas nécessaire lorsque le catalogue ou l'ANStat fournit directement
la livraison harmonisée. Le pipeline R commence après cette harmonisation et
ne doit pas être lancé sur des fichiers bruts simplement renommés.

\subsection{Architecture et inventaire du pipeline}

L'ensemble des résultats est produit par un pipeline R de vingt-six étapes. Les
données brutes restent dans \path{01_data_sources}; les transformations
intermédiaires sont écrites en parquet dans \path{02_data_intermediate}; les
tableaux et figures sont déposés dans \path{07_reports}. Le fichier
\path{05_scripts_R/00_setup.R} centralise les chemins, charge les dépendances et
les fonctions communes. Le fichier \path{05_scripts_R/00_master.R} déclare
l'ordre d'exécution et expose l'orchestrateur \texttt{lance\_pipeline}. Son
chargement prépare l'environnement sans lancer de simulation.

\begin{longtable}{>{\centering\arraybackslash}p{1.2cm}
  >{\RaggedRight\arraybackslash}p{4.5cm}
  >{\RaggedRight\arraybackslash}p{8.0cm}}
\caption{Inventaire des étapes du pipeline R}
\label{tab:inventaire-pipeline}\\
\toprule
Étape & Script principal & Objet \\
\midrule
\endfirsthead
\toprule
Étape & Script principal & Objet \\
\midrule
\endhead
1 & \path{01_prepare_data.R} & Import, contrôles, achats observés et
winsorisation de la consommation. \\
2 & \path{02_mapping_tax.R} & Appariement des postes de consommation aux taux
statutaires de TVA. \\
3 & \path{03_compute_taxes.R} & Calcul de la TVA directe au niveau du ménage. \\
4 & \path{04_analysis.R} & Profils par décile, quintile, milieu et région. \\
5 & \path{05_progressivity.R} & Gini, concentration, Kakwani et
Reynolds--Smolensky. \\
6 & \path{06_01_sensitivity_taxation.R} &
Scénarios de taxation effective liés à l'informalité des achats. \\
7 & \path{06_02_sensitivity_ranking.R} &
Sensibilité aux unités et variables de classement. \\
8 & \path{07_appendix_tables.R} & Tableaux distributifs complémentaires. \\
9 & \path{08_figures.R} & Figures principales de l'incidence de la TVA. \\
10 & \path{09_vat_determinants.R} & Corrélats du taux effectif de TVA. \\
11 & \path{10_reform_chicken_inputs.R} & Simulation sur les intrants avicoles. \\
12 & \path{11_reform_figures.R} & Figures associées à cette simulation. \\
13 & \path{12_poverty_incidence.R} & Indices FGT avant et après fiscalité
indirecte. \\
14 & \path{13_leontief_io.R}, \path{13b_leontief_local_io.R} &
Propagation de la TVA non déductible avec ICIO et avec les TRE ivoiriens. \\
15 & \path{14_leontief_poverty.R} & Pauvreté après incorporation de la TVA
enchâssée. \\
16 & \path{15_reform_vat_simulation.R} & Taxe de 9~\% sur la vente finale à
TVA incorporée constante, scénarios d'informalité et recyclage budgétairement
neutre. \\
17 & \path{16_direct_taxes.R} & Impôts directs, cotisations sociales et
diagnostics d'assujettissement. \\
18 & \path{17_indirect_other.R} & Accises, droits de douane, cascade TTC et
validation des concordances TRE--TEC. \\
19 & \path{18_transfers.R} & Pensions, transferts directs, filets sociaux et
reconstruction des branches PDI et PGT. \\
20 & \path{19_subsidies.R} & Réductions de prix de l'électricité et de l'eau;
carburants en robustesse seulement. \\
21 & \path{20_inkind_education.R} & Dépenses publiques d'éducation attribuées
par niveau, brutes et nettes des paiements. \\
22 & \path{21_inkind_health.R} & Soins ambulatoires moyens, hospitalisations
observées, paiements directs et robustesses. \\
23 & \path{22_income_concepts.R} & Assemblage et contrôle des revenus primaire,
net, brut, disponible, consommable et final. \\
24 & \path{23_marginal_contribution.R} & Contributions avec/sans, concentration
et décomposition exacte de Shapley. \\
25 & \path{24_fiscal_impoverishment.R} & Appauvrissement fiscal, gains sous le
seuil, dominance et inférence Rao--Wu. \\
26 & \path{25_ceq_report_tables.R} & Classeur CEQ final, contrôle macroéconomique
ANStat, dictionnaire et manifeste de réplication. \\
\bottomrule
\end{longtable}

Les fonctions distributives partagées se trouvent dans
\path{05_scripts_R/utils/distributive.R}; la résolution du modèle de prix, les
scénarios de TVA et la concordance TRE sont isolés dans
\path{utils/leontief_vat.R}, \path{utils/vat_scenarios.R} et
\path{utils/tre_mapping.R}. Les barèmes de 2021 sont documentés séparément
dans \path{00_documentation/methodologie/MEMOIRE_FISCAL_CIV2021.md}. Cette
organisation permet de relier chaque résultat à une étape tout en évitant de
réimprimer plusieurs centaines de lignes de code dans le papier.

\subsection{Installation, contrôles et exécution}

L'environnement de référence utilise R~4.5.3. Le fichier \path{renv.lock} fixe
les versions des paquets et \path{replication_package/environment/sessionInfo.txt}
documente la plateforme de référence. Depuis la racine du dépôt, un économiste
peut restaurer cet environnement puis lancer les contrôles par :

\begin{verbatim}
Rscript --vanilla replication_package/code/00_restore_environment.R
Rscript --vanilla replication_package/code/00_preflight.R
\end{verbatim}

Le premier programme installe les versions enregistrées dans la bibliothèque
isolée du projet. Le second vérifie la version de R, les paquets, l'arborescence,
les vingt entrées attendues, leur taille et leur empreinte MD5. Il compare ces
informations au manifeste \path{replication_package/data/data_manifest.csv} et
s'arrête dès qu'un contrôle obligatoire échoue. Les microdonnées ne sont jamais
recopiées dans le paquet; le fichier
\path{replication_package/data/access-restricted-data.md} explique leur accès et
leur emplacement.

Après un pré-contrôle sans échec, une seule commande reconstruit les vingt-six
étapes :

\begin{verbatim}
Rscript --vanilla replication_package/code/00_run_all.R
\end{verbatim}

Ce point d'entrée extrait automatiquement les matrices des classeurs TRE si les
fichiers intermédiaires n'existent pas, exécute les étapes dans leur ordre,
conserve le journal et appelle la vérification finale. Il n'installe aucun paquet
silencieusement. Pour contrôler des sorties déjà produites sans relancer les
calculs, la commande est :

\begin{verbatim}
Rscript --vanilla replication_package/code/01_verify_outputs.R
\end{verbatim}

Le vérificateur compare les indicateurs centraux aux valeurs attendues dans
\path{replication_package/output/expected_metrics.csv}. Il contrôle notamment
les Gini, les taux de pauvreté, l'appauvrissement fiscal, la masse du PSSN et la
décomposition de Shapley. Les rapports
\path{preflight_report.csv} et \path{verification_report.csv} doivent contenir
zéro échec. Une exécution d'une plage, par exemple
\texttt{lance\_pipeline(17,17)}, reste utile pour développer un module, mais ne
constitue pas une réplication depuis les entrées.

\subsection{Deux parcours de réplication : lecteur et agent}

Le fichier \path{replication_package/README.md} donne à un économiste la liste
des données, les commandes, les sorties attendues et le lien entre chaque
tableau ou figure et son programme. Le manifeste
\path{replication_package/exhibit_map.csv} fournit la même correspondance sous
forme tabulaire. Les besoins en mémoire, espace disque, logiciels et temps
d'exécution sont consignés dans
\path{replication_package/environment/computational_requirements.md}.

Un agent de programmation suit exactement le même point d'entrée. Le fichier
\path{AGENTS.md} fixe les règles du dépôt; le guide
\path{replication_package/AGENT_RUNBOOK.md} décrit les contrôles autorisés, les
données confidentielles et les conditions d'arrêt; enfin,
\path{replication_package/replication_spec.json} expose les entrées, commandes
et critères de succès dans un format lisible par machine. Ces fichiers ne créent
pas une seconde méthode : ils traduisent le protocole humain sans modifier les
hypothèses ni les tolérances numériques.

Le manuscrit est reconstruit depuis
\path{00_documentation/working_paper/} par un passage \texttt{pdflatex}, un
passage \texttt{bibtex}, puis deux nouveaux passages \texttt{pdflatex}. La
révision Git, le journal de calcul, \texttt{sessionInfo()} et les deux rapports
de contrôle doivent accompagner les résultats diffusés. Une réplication publique
complète exige enfin que la version utilisée soit gelée dans une release et que
la licence de diffusion du code soit explicitement approuvée.
\end{document}
